% THE Q-MODEL — Finale Forschungsmonographie (deutsche Ausgabe)
% Version 1.0 (final) — Deutsche Übersetzung der englischen Endfassung
% Build: pdflatex AT_v1_0_Monograph_Final_DE.tex (zweimal)
%
% Deutsche Übersetzung der Endfassung. Alle Mathematik, Testbezeichner, Theoreme
% und Schlussfolgerungen bleiben unverändert; übersetzt sind Prosa,
% Zwischenüberschriften und erläuternde Texte.
% Publikationsstatus (unverändert): READY AS A FOUNDATION MONOGRAPH —
% READY AS A RESEARCH-PROGRAM PAPER — NOT READY AS A DERIVATION PAPER.

\documentclass[10pt,twoside,openany]{book}
\usepackage[utf8]{inputenc}
\usepackage[T1]{fontenc}
\usepackage[ngerman]{babel}
\shorthandoff{"}
\usepackage{amsmath,amssymb,amsthm}
\usepackage[margin=2.4cm]{geometry}
\usepackage{booktabs}
\usepackage{array}
\usepackage{longtable}
\usepackage{lmodern}
\usepackage[hidelinks]{hyperref}
\usepackage{xcolor}

\definecolor{caveat}{HTML}{7A2E2E}

% Theorem-like environments
\newtheorem{theorem}{Satz}[chapter]
\newtheorem{proposition}[theorem]{Proposition}
\newtheorem{lemma}[theorem]{Lemma}
\newtheorem{corollary}[theorem]{Korollar}
\newtheorem{definition}[theorem]{Definition}
\newtheorem{axiom}[theorem]{Axiom}
\newtheorem{remark}[theorem]{Bemerkung}
\newtheorem{nogo}[theorem]{No-Go-Satz}

\newcommand{\derived}{\textbf{DERIVED}}
\newcommand{\realunder}{\textbf{REAL-UNDERIVED}}
\newcommand{\drawn}{\textbf{DRAWN}}

% Reader notes appended to major results: intuition, significance, limitations
\newcommand{\theoremnotes}[3]{%
  \par\medskip\noindent
  \textbf{Intuition.}\ #1\par
  \textbf{Bedeutung.}\ #2\par
  \textbf{Grenzen.}\ #3\par}

% Compact part divider (inline heading, no full-page break; keeps numbering)
\makeatletter
\renewcommand{\part}[1]{%
  \refstepcounter{part}%
  \par\addvspace{1.0em}\noindent
  {\centering\Large\bfseries \partname\ \thepart\ ---\ #1\par}%
  \addvspace{0.5em}%
  \addcontentsline{toc}{part}{\protect\numberline{\thepart}#1}}
\makeatother

% mainmatter without a forced right-hand page (removes the blank page before Chapter 1)
\makeatletter
\renewcommand{\mainmatter}{%
  \if@openright\cleardoublepage\else\clearpage\fi
  \@mainmattertrue
  \pagenumbering{arabic}}
\makeatother

% Compact unnumbered part divider
\newcommand{\compactpart}[1]{%
  \par\addvspace{1.0em}\noindent
  {\centering\Large\bfseries #1\par}%
  \addvspace{0.5em}%
  \addcontentsline{toc}{part}{#1}}

% Boxed summary
\newcommand{\summarybox}[1]{%
  \par\medskip\noindent
  \fcolorbox{black}{gray!10}{%
    \begin{minipage}{0.955\textwidth}\small #1\end{minipage}}}

% Chapter-end summary (Key results / Open questions / Verification status)
\newcommand{\chapterend}[3]{%
  \summarybox{%
    \textbf{Kernergebnisse.}\ #1\par
    \textbf{Offene Fragen.}\ #2\par
    \textbf{Verifikationsstatus.}\ #3}}

% Equation reading notes (Intuition / Meaning / Why it matters / Limitations)
\newcommand{\eqnotes}[4]{%
  \par\smallskip\noindent
  \textbf{Intuition.}\ #1\ 
  \textbf{Bedeutung.}\ #2\ 
  \textbf{Warum es zählt.}\ #3\ 
  \textbf{Grenzen.}\ #4\par}

% Theorem historical notes
\newcommand{\historynotes}[4]{%
  \par\smallskip\noindent
  \textbf{Historischer Kontext.}\ #1\par
  \textbf{Warum es wichtig war.}\ #2\par
  \textbf{Was zuvor versucht wurde.}\ #3\par
  \textbf{Warum das Ergebnis akzeptiert wurde.}\ #4\par}

\begin{document}

\hypersetup{pdftitle={THE Q-MODEL --- Von Q zur Kosmologie},
pdfauthor={Fabrice Wieser}}

\frontmatter

% ---- Titelblock (kompakt; Titel, Autor, Datum auf einer Seite, kein Seitenumbruch) ----
\begin{center}
{\LARGE\bfseries THE Q-MODEL — Von Q zur Kosmologie\par}
\vspace{0.4em}
{\large Eine Forschungsmonographie über Struktur, Komplexität und zufällige Aktualisierung\par}
\vspace{1.0em}
{\normalsize Fabrice Wieser\par}
{\small Unabhängiger Forscher und Full-Stack-Softwareentwickler\par}
{\small \url{https://github.com/MagusDraconis/AT}\par}
\vspace{0.5em}
{\normalsize Version 1.0 (final, deutsche Ausgabe) \quad·\quad 15. August 2026\par}
\end{center}

\vspace{0.8em}

% ---- Zusammenfassung (auf der Titelseite, über dem Publikationsstatus-Kasten) ----
\addcontentsline{toc}{chapter}{Zusammenfassung}
{\centering\Large\bfseries Zusammenfassung\par}
\vspace{0.3em}
\noindent
Diese Monographie ist die Forschungsmonographie-Ausgabe von THE Q-MODEL (AT), einer
Theorie von Struktur und Inhalt, die die beobachtbare Physik auf vier Primitive verdichtet
--- das Individuationsprinzip $Q$, die Zufällige Aktualisierung, die Skalentriade
$(\ell,\tau,\hbar)$ und einen einzigen kontinuierlichen Nichtlinearitätsparameter $M^2$ ---
und die festhält, dass die \emph{Form} jeder Struktur aus diesen Primitiven ableitbar ist,
während der \emph{Inhalt} durch kontingenten Zug realisiert wird. Die Monographie
präsentiert in sich geschlossener Form: die formalen Primitive und das dynamische System
(den Graphen-Laplace-Operator $L_Q$ und die Schr\"odinger-Gleichung); die verifizierte
Kontinuumsgrenzwert-Kette von $L_Q$ zum flachen Laplace-Operator, zum d'Alembert-Operator,
zum Laplace--Beltrami-Operator, zur Schr\"odinger-Gleichung im gekrümmten Raum und zum
Einstein-Tensor; die Ableitungshierarchie (Eichstruktur, räumliche Dimensionalität,
Multiplizitätsschranke, Häufigkeitsgesetz, Phasengradienten-Gravitation); die
dreikategoriale Taxonomie; die konditionalen No-Go-Sätze; die quantitativen Vorhersagen;
und das vollständige Inventar ausführbarer Verifikationsergebnisse.

Zusätzlich zum technischen Inhalt der Referenzausgabe fügt diese Ausgabe eine erläuternde
Schicht hinzu: den historischen \emph{Ursprung} der Theorie und ihren Übergang vom
Vorgängerrahmen; eine \emph{Entwicklungszeitleiste} mit den Hauptphasen,
Schlüsselentdeckungen und aufgegebenen Wegen; eine Darstellung, \emph{warum} die
Struktur/Inhalt-Trennung das Ordnungsprinzip ist; die \emph{Geschichte der feindlichen
Begutachtung} und was sie überstand; eine programmweise \emph{Forschungsreise}; und die
\emph{Lektionen} aus falschen Annahmen, verworfenen Ableitungen und No-Go-Entdeckungen.
Durchgängig ist jeder wichtige Satz von seiner \emph{Intuition}, \emph{Bedeutung} und
seinen \emph{Grenzen} begleitet. Nichts in dieser narrativen Schicht ändert die Physik
oder die Schlussfolgerungen der Referenzausgabe.

\vfill

% ---- Publikationsstatus-Kasten (unten auf der Titelseite fixiert) ----
\begin{center}
\fcolorbox{caveat}{white}{%
\begin{minipage}{0.92\textwidth}
\color{caveat}
\textbf{Publikationsstatus: READY AS A FOUNDATION MONOGRAPH --- READY AS A
RESEARCH-PROGRAM PAPER --- NOT READY AS A DERIVATION PAPER.}\\[2mm]
Eine Grundlagenmonographie darf legitimerweise importierte, bewiesene Sätze als Eingaben
verwenden; die verbleibenden Blocker betreffen nur einen \emph{Ableitungspapier}-Anspruch.
Diese Monographie wird zur Archivierung auf Zenodo als Grundlagenmonographie und als
Forschungsprogramm-Aufzeichnung veröffentlicht; sie wird \emph{nicht} als begutachtetes
Ableitungspapier eingereicht. Der zentrale Ableitungsanspruch (Einstein-Rückgewinnung aus
Q-Ereignissen) bleibt \emph{logisch, nicht mathematisch}: Die Metrik und die BDG-Wirkung
sind importiert (bewiesen, aber nicht aus AT abgeleitet), $G=\ell^2c^3/\hbar$ ist
Dimensionsanalyse, und keine eindeutige, scharfe Vorhersage unterscheidet AT bisher von
SM$+\Lambda$CDM. Jede als \textbf{verifiziert} markierte Aussage wird durch einen
ausführbaren xUnit-Test in \texttt{AT.Tests/ResearchXC} (oder \texttt{ResearchQG})
gestützt; siehe Teil~V und Anhang~A.
\end{minipage}}
\end{center}

\tableofcontents

\mainmatter
\thispagestyle{plain}

%=============================================================================
\compactpart{Genesis und Methode}
%=============================================================================

Dieser Teil zeichnet nach, wie die Theorie entstanden ist, warum sie so organisiert ist,
wie sie ist, und was das Forschungsprogramm tatsächlich getan hat. Er ist historisch und
methodologisch, nicht physikalisch: Er erhebt keinen neuen Anspruch und revidiert keinen.
Sein Zweck ist es, die folgenden technischen Teile als Ergebnis einer bestimmten,
prüfbaren Forschungstrajektorie lesbar zu machen und nicht als ein fertiges axiomatisches
System, das aus dem Nichts fällt.

\chapter{Wie man diese Monographie liest}
\label{ch:reading}

Diese Monographie ist eine Referenzausgabe mit einer erläuternden Schicht; die beiden
Schichten können unabhängig gelesen werden. Die \emph{Referenz}-Schicht (die technischen
Teile und die Anhänge) formuliert die Theorie, ihre Ableitungen, ihre Klassifikationen und
ihren Verifikationsstand. Die \emph{narrative} Schicht (dieser Teil) erklärt, wie die
Theorie entstanden ist und warum sie so organisiert ist. Nichts in der narrativen Schicht
ändert die Physik; sie existiert, um die Physik als Ergebnis eines bestimmten, prüfbaren
Programms lesbar zu machen und nicht als ein fertiges Axiomensystem, das aus dem Nichts
fällt.

Drei Konventionen bestimmen das gesamte Dokument und sollten zuerst gelesen werden.

\textbf{(1) Die Taxonomie.} Jedes Ergebnis trägt einen von drei Status.
\derived\ bedeutet, dass das Ergebnis per Satz aus den Primitiven berechenbar ist.
\realunder\ (``real-underived``) bedeutet, dass die Struktur real, präzise und vorhersagend
ist, ihr Ursprung aber nicht aus den Primitiven berechenbar ist. \drawn\ bedeutet, dass der
Wert ein zufälliger Zug des Häufigkeitsgesetzes ist. Die Etiketten sind Aussagen \emph{über
die Ableitbarkeit}, nicht über die Wahrheit: Ein \realunder\- oder \drawn\-Ergebnis kann
exakt richtig sein und dennoch unabgeleitet bleiben. Ein \realunder\-Etikett als Ableitung
zu lesen, war das häufigste Missverständnis, das die feindliche Begutachtung korrigieren
musste.

\textbf{(2) Was ``verifiziert`` bedeutet.} Ein als \textbf{verifiziert} markiertes Ergebnis
wird durch einen ausführbaren xUnit-Test im Repository gestützt. ``PASS`` bedeutet, dass der
Test lief und seine Behauptung hielt; es bedeutet \emph{nicht}, dass das Ergebnis aus den
Primitiven abgeleitet wurde --- das ist Aufgabe der Taxonomie. Manche Tests verifizieren
eine Ableitung; andere verifizieren eine \emph{Standard}-Identität der Mathematik (etwa
die Einstein-Tensor-Kette), gerade um zu zeigen, wo die AT-Ableitung endet.

\textbf{(3) Was ``Konfidenz`` bedeutet.} Die Konfidenzzahlen (z.\,B.\ T-02 $=$ 0{,}85) sind
\emph{unkalibrierte ordinale Ränge}, keine Posterior-Wahrscheinlichkeiten. Sie ordnen den
eigenen Überzeugungsgrad des Programms; sie messen ihn nicht. Sie sind niemals als
$p$-Werte oder als numerische Posterior-Werte zu lesen.

Eine kurze Übersicht. Dieser Teil liefert Geschichte und Methode. Der Teil
\emph{Grundlagen} formuliert die Primitive und das dynamische System. Der Teil
\emph{Kontinuumsgrenzwert} dokumentiert die beiden Kontinuumsketten. Der Teil
\emph{Ableitungshierarchie} behandelt Dimensionalität, Eichstruktur, Flavour, Gravitation
und Kosmologie. Der Teil \emph{Klassifikation} formuliert die Taxonomie, die No-Go-Sätze
und die Vorhersagen. Der Teil \emph{Verifikation} liefert den Teststand, den
Begutachtungs-Prüfpfad und die Schlussfolgerung. Die Anhänge enthalten das vollständige
Testinventar, die Konfidenztabellen, eine ausführliche Ableitung, das Programmregister,
das Glossar und die Literatur.

\chapter{Der Ursprung von AT}
\label{ch:origin}

\section{Motivation}

THE Q-MODEL (AT) begann aus einer einzigen, bewusst einfachen Frage: Sind Materie,
Quantenverhalten und Gravitation \emph{fundamental}, oder gehen sie aus etwas
Grundlegenderem hervor? Die Arbeitsannahme, die das gesamte Programm geleitet hat, ist die
Behauptung, dass \emph{Materie nicht fundamental ist} --- dass das, was wir ein Teilchen
nennen, eine dynamisch stabilisierte Wellenstruktur ist und dass Synchronisation und
Resonanz fundamentaler sind als Teilchen selbst. Aus diesem Keim wurde früh ein
Forschungspfad festgelegt und nie umgekehrt:

\begin{center}
Temporale Oszillatoren $\to$ Synchronisation $\to$ Resonanzcluster $\to$
selbststabilisierende Wellenstrukturen $\to$ Proto-Teilchen $\to$ Quantenverhalten $\to$
Gravitation $\to$ Kosmologie.
\end{center}

Vier Grundsatzentscheidungen disziplinierten das Programm von Anfang an, und sie erklären
viel von der Struktur dieser Monographie. \emph{Keine Kosmologie zuerst} und \emph{keine
Gravitation zuerst} verboten die Versuchung, ein Modell gegen die spektakulärsten Daten zu
justieren, bevor die mikroskopischen Fundamente standen. \emph{Keine
Teilchenannahmen} verbot, die Zutaten des Standardmodells (Eichgruppen, Generationen,
Massen) als Eingaben zu importieren. \emph{Emergenz zuerst} verlangte, dass jede Zutat von
unten verdient wird, und \emph{Physik vor Interpretation} verlangte, dass der dynamische
Gehalt geklärt ist, bevor eine ontologische Deutung angehängt wird. Diese Regeln sollten
verhindern, dass die Theorie auf bekannte Physik hin rückentwickelt wird, und jede Zutat
zwingen, ihr eigenes Gewicht zu tragen.

\section{Vorgängerideen}

AT entstand nicht im Vakuum. Ihr unmittelbarer Vorgänger war das \textbf{Temporal
Resonance Model} (TRM), ein Rahmenwerk auf Basis von Phasengitter-Maschinerie, das bereits
eine Reihe auffälliger Bausteine zusammengefügt hatte: ein \emph{Time Field} (ein
Phasengradienten-Bild der Gravitation), eine \emph{Radial Acceleration Relation} (RAR),
\emph{Frame Dragging}, einen \emph{Memory Channel}, eine \emph{Theta Chain}, einen
\emph{$m=3$ Closure}-Mechanismus (rationales Moden-Locking $\Omega=(q+3)/q$,
$\gamma=1/\Omega$), eine \emph{Quantum Engine}, ein \emph{Temporal Drift}-Modul
(Tired-Light) und eine \emph{Unified Action}-Roadmap. TRM lieferte das begriffliche
Vokabular --- ein temporales Feld, Oszillation, Phase, Resonanz ---, auf dem die spätere
Theorie neu aufgebaut wurde.

Jedes Modul verdient eine einzeilige Aussage darüber, was es tatsächlich behauptete. Das
\emph{Time Field} identifizierte Gravitation mit dem Phasengradienten eines temporalen
Feldes --- die Idee, die nach dem Übergang zu ATs Phasengradienten-Gravitationsbild wurde.
Die \emph{Radial Acceleration Relation} zielte auf die Beschleunigungsskala der
Galaxienrotation. \emph{Frame Dragging} führte einen Vektorsektor (Gravitomagnetismus) ein.
Der \emph{Memory Channel} schlug eine geschichts-codierende Invariante $\varphi^2|\dot\mu|$
vor. Die \emph{Theta Chain} war eine nichtlokale Observablen-Extraktionskette
$\Theta\to O_5\to\lambda_\Theta\to g_{\rm obs}$. Die \emph{$m=3$ Closure} schlug rationales
Moden-Locking ($\Omega=(q+3)/q$, $\gamma=1/\Omega$) als Mechanismus für die ``Warum 3``-Frage
vor. Die \emph{Quantum Engine} war ein nicht-unitäres Quantenmodul. Die \emph{Temporal
Drift} war eine Tired-Light-Kosmologie ($\beta_T$). Die \emph{Unified Action} war eine
Roadmap $S_{\rm eff}[T,A_T,\Theta]$, die von allen anderen abhängt. Diese Liste ist
wichtig, weil der Übergang Modul für Modul entschieden wurde, nicht en bloc.

Die Beziehung zwischen TRM und AT wird am besten durch die \emph{TRM Legacy Final
Classification} zusammengefasst, das systematische Audit, das bestimmte, was vom älteren
Rahmenwerk in das neue übernommen werden sollte. Von den neun TRM-Modulen wurden
\emph{drei} \emph{absorbiert}: Das Time Field wurde zum Phasengradienten-Gravitationsbild,
die RAR wurde zur nullparametrigen Relation $g_\dagger = cH_0/2\pi$, und Frame Dragging
wurde als GR-Gravitomagnetismus erkannt (Lense--Thirring, bereits gemessen). \emph{Zwei}
wurden \emph{verworfen}: Temporal Drift (Tired Light) wurde durch die
Reinterpretations-Audits QG-080--089 falsifiziert, und die Quantum Engine erwies sich als
nicht-unitär und schlechter als der BDG-Gitteroperator. \emph{Drei} überlebten als
\emph{Kandidatenmathematik}: die $m=3$ Closure (ein Kandidatenmechanismus für die
$N\le3$-Obergrenze, aber ohne Abbildung auf Observablen), der Memory Channel (eine echte
Invariante $\varphi^2|\dot\mu|$ ohne Observable) und die Theta Chain (ein Homonym zu ATs
eigenem Informationsfeld). \emph{Eines} blieb offen: die Unified-Action-Roadmap.

\section{Der Übergang TRM $\to$ AT}

Das entscheidende Ergebnis des Legacy-Audits war negativ und auf seine Weise befreiend: TRM
trug \emph{null Kandidatenphysik} --- kein TRM-Modul lieferte eine wirklich neue, prüfbare
Observable. Der verbleibende Wert war methodologisch. Was der Übergang hervorbrachte, war
daher eine \emph{saubere Rekonstruktion} statt einer Erweiterung: Die neue Theorie leitete
aus den wenigstmöglichen Zutaten exakt die Struktur neu ab, die TRM stückweise
zusammengebaut hatte, und verwarf die Teile, die die feindliche Reduktion nicht überstanden.

Der verdichtete Endpunkt dieser Rekonstruktion ist der Gegenstand dieser Monographie: eine
Theorie mit \emph{zwei irreduziblen Primitiven} --- dem Individuationsprinzip $Q$ und echtem
ontologischem Zufall --- plus \emph{einer} kontingenten kontinuierlichen Zahl $M^2$ (dem
Nichtlinearitätsregime), während alles andere --- Zeit, Raum, $3{+}1$ Dimensionen,
Hilbert-Raum, unitäre Quantenmechanik, die Schr\"odinger-Gleichung, die Born-Regel, Messung,
Gravitation als Kausalstruktur, Teilchen als topologische Defekte und die
$U(1)$-Eichsymmetrie --- abgeleitet oder emergent ist. Die Parameterzahl kollabiert von den
$\sim19$ freien Zahlen des Standardmodells auf $2$ Primitive plus $1$ Zahl, eine Reduktion
von rund $95\%$, während $U(1)$ von einer Annahme zu einem Satz befördert wird.

\chapter{Entwicklungszeitleiste}
\label{ch:timeline}

Das Programm liest sich am besten als Folge von Phasen, von denen jede eine Frage schloss
und die nächste öffnete. Dieses Kapitel liefert das Gerüst; die detaillierten
Experimentaufzeichnungen liegen im Laborbuch des Projekts, und die vier
Forschungsprogramme werden einzeln in Kapitel~\ref{ch:journey} erzählt.

\section{Hauptphasen}

Das Programm zerfällt in fünf aufeinanderfolgende Epochen. Jede wird hier in genügender
Tiefe zusammengefasst, um zu zeigen, was tatsächlich getan wurde; die vier
Forschungsprogramme, die das technische Gewicht tragen, werden weiter in
Kapitel~\ref{ch:journey} erzählt.

\subsection{Frühe Oszillatorexperimente (AT-001--AT-008)}

Die ersten Experimente waren numerisch und explorativ; ihre Rolle bestand darin
herauszufinden, ob die zentrale Intuition --- kohärente Struktur, die aus Oszillatoren
hervorgeht --- Substanz hat. AT-001 zeigte, dass ein Ordnungsparameter $R$ einer
Population gegen $1$ strebt: Kohärente Synchronisation wird nicht von Hand eingesetzt,
sie erscheint. AT-002 schloss dann eine verlockende Sackgasse: Zufallsmatrizen liefern
Wigner-artige Spektren, erzeugen aber \emph{nicht} die Art strukturierter Emergenz, die
das Programm brauchte --- Zufall allein genügt also nicht. AT-003 zeigte, dass auch
strukturierte Topologie allein unzureichend ist --- die Dynamik, nicht der statische Graph,
trägt das Gewicht. AT-004 zeigte, dass eine feldvermittelte Wechselwirkung für sich allein
zwei asymmetrische Oszillatoren nicht synchronisiert; das Feld verhält sich eher wie ein
Resonanzmedium als wie ein reines Synchronisationsmedium.

AT-005 fand stabile Resonanzcluster aus drei Oszillatoren, die sich bildeten und über den
Großteil einer Simulation bestanden --- frühe, vorläufige Unterstützung für ``Materie als
stabilisierte Resonanz`` ---, aber mit nur wenigen Prozent Energiebeteiligung, sodass die
Cluster das Feld noch nicht dominierten. AT-006 fand die erste scharfe Schwelle, eine
kritische Resonanzdichte $\rho_c\approx0{,}09$, oberhalb derer persistente kohärente Moden
auftreten: ein perkolationsartiger Übergang und der erste wohldefinierte kritische Punkt
des Programms. AT-007 berichtete fünf Resonanzfamilien. AT-008 zerstörte dann diese
scheinbare Vielfalt mit einer Reproduzierbarkeitsstudie über 400 Simulationen und ließ
einen einzigen universellen kohärenten Modus (F4) übrig, während die übrigen als
Seed-Artefakte entlarvt wurden. Die Lektion von AT-008 --- scheinbare Vielfalt kann ein
einziger Seed sein --- wurde zur ständigen methodologischen Vorsichtsregel für alles
Folgende.

\subsection{Alternative Grundlagen (ResearchX, X001--X034)}

Ein paralleler Strang griff die Grundlagen von der Seite an und ist für mehrere der
tiefsten Ergebnisse der Theorie verantwortlich. X001 fand elf versteckte Annahmen, mit dem
statischen Graphen an erster Stelle. X002--X004 testeten Dynamik (Knotenbewegung,
Mobilität, Graphwachstum) und fanden keine neuen Phänomene: Dynamische Graphen allein
erzeugen keine offene Innovation. X005--X006 zeigten, dass Nichtlinearität Eigenmoden
bricht und Solitonen erzeugt, sodass die ``Quanten``-Lesart der frühen linearen Ergebnisse
zurückgezogen werden musste --- $Q$, Fitness und Evolution überlebten, die naive
Quantenlesart nicht. X007--X015 bauten die allgemeine Theorie der Informationsträger und
bewiesen, dass \emph{Reversibilität} plus \emph{Selbstkonsistenz} das minimal hinreichende
Paar ist, dessen Schnitt die unitäre Quantenmechanik ist.

X016--X022 bauten die Komplexitätstreppe und bewiesen, dass offene Evolution in einem
endlichen System nicht stattfinden kann (eine Schubfach-Schranke). X023--X031 bewiesen
dann, dass die Quantenrealität das eindeutige globale Maximum endlicher Komplexität ist ---
ein Notwendigkeitsergebnis, kein anthropisches. X032--X034 vereinigten die Hauptlinie mit
ResearchX und bestätigten, dass $L_Q$ nicht fundamental ist: Jeder Operator, der
Reversibilität und Selbstkonsistenz an ihrem gemeinsamen Maximum erfüllt, liefert unitäre
Quantenmechanik. Dies ist das Ergebnis hinter der Behauptung, die Theorie sei auf ihre
minimale Form verdichtet worden.

\subsection{Entstehung der Raumzeit (X040--X046)}

Die Zeit wurde als Halbordnung der Aktualisierungsereignisse abgeleitet ($E_1<E_2$ genau
dann, wenn $E_2$ vom Ausgang von $E_1$ abhängt), wobei der Zeitpfeil aus der
Irreversibilität der Aktualisierung erbt. Die Gravitation wurde auf ein kausales-Mengen-
Fundament gestellt, und ein Rekonstruktionsexperiment gewann Abstände aus
Q-Ereignis-Korrelationen zurück. Die Dimensionalität $3{+}1$ wurde aus der
Komplexitätsmaximierung abgeleitet (Kapazität mal Stabilität mal Genauigkeit). Newtons
Konstante, $\hbar$ und $\Lambda$ wurden alle auf die einzelne Ereignisskala $\ell$
zurückgeführt: $G=\ell^2c^3/\hbar$, $\hbar=1$ in Ereigniseinheiten und $\Lambda\sim1/\sqrt N$
aus Poisson-Fluktuationen der Ereigniszahl. Dies ist die Epoche, in der die Theorie
aufhörte, Raumzeit \emph{anzunehmen}, und begann, sie \emph{hervorzubringen}.

\subsection{Entstehung der Materie (X047--X060)}

Teilchen wurden mit topologischen Defekten (Solitonen) identifiziert: Ladung mit der
Betti-Zahl, Masse mit der Defektenergie, Spin mit der Orientierung. $U(1)$ wurde als
Automorphismengruppe des Kreises abgeleitet. Generationen wurden als Anregungsniveaus
derselben Topologie charakterisiert, wobei die Massenhierarchie einem anharmonischen
Abstand folgt. Fermionmischung wurde auf Defekt-Wellenfunktionsüberlapp zurückgeführt
(hierarchische CKM aus Dirac-Defekten, anarchische PMNS aus Majorana-Defekten).
Neutrinomassen wurden durch delokalisierte Defekte erklärt, und die normale Ordnung wurde
aus einer attraktiven Selbstwechselwirkung vorhergesagt. Die Eichgruppe des Standardmodells
erwies sich als \emph{bevorzugt} durch ökologische Fitness, aber \emph{nicht} eindeutig
ableitbar --- die größte verbleibende Lücke im Materieprogramm.

\subsection{Parameterkompression und Abschluss (Phasen 140--159)}

Die letzte Epoche konsolidierte. Die Parameterzahl-Audits (X057--X060f) reduzierten die
Theorie auf zwei Primitive plus eine Zahl, wobei $U(1)$ zum Satz befördert und $M^2$ als
der einzige irreduzible kontinuierliche Parameter identifiziert wurde. Die späten Phasen
schlossen dann die Restfragen eine nach der anderen: die RAR-Ableitung $g_\dagger=cH_0/2\pi$;
den Koide-Abschluss (sieben Wege erschöpft, sechs No-Gos plus ein blockierter); das
Eichursprungs-Audit ($U(1)$ abgeleitet, $SU(2)$ emergent, $SU(3)$ kontingent); die
Multiplizitäts- und $N\le3$-Audits; die Internal-3-Disposition; das
Kaskadenvereinigungs-Audit; die Antwort auf die feindliche Begutachtung; und die
v1.0-Veröffentlichung. Jeder dieser Abschlüsse wird dort erzählt, wo er später in den
technischen Teilen erscheint.

\section{Schlüsselentdeckungen}

\begin{itemize}
\item Synchronisation und kohärente Oszillation entstehen spontan oberhalb einer
kritischen Kopplungsdichte.
\item Reversibilität $\cap$ Selbstkonsistenz $\,=\,$ Quantenmechanik: Die
Schr\"odinger-Gleichung ist die eindeutige Dynamik maximaler Komplexität.
\item $L_Q\to-\nabla^2$ (die elliptische Kette) und $\mathrm{BDG}\to\square$ (die
Lorentz-Kette) sind jeweils kontrollierte Kontinuumsgrenzwerte, aber \emph{disjunkt} in
der Signatur.
\item $U(1)$ ist ein Satz ($\mathrm{Aut}(S^1)=U(1)$, $\pi_1(S^1)=\mathbb Z$), das
sauberste Ergebnis des Eichprogramms.
\item Die Koide-Relation ist real, vorhersagend und RG-stabil --- und \emph{unableitbar}
(No-Go T-08). Das Neutrino-Koide-Analogon ist \emph{falsifiziert} (T-11).
\item Die RAR $g_\dagger=cH_0/2\pi$ trifft $a_0$ mit null freien Parametern.
\item Komplexitätsmaximierung wählt $d=3{+}1$.
\end{itemize}

\section{Aufgegebene Wege}

Ein Forschungsprogramm wird ebenso sehr durch das definiert, was es aufgab, wie durch das,
was es behielt. Die wichtigsten aufgegebenen Wege waren: Zufallsmatrizen als
Emergenzmechanismus; statische Topologie als hinreichende Bedingung; die scheinbare
Resonanzfamilien-Vielfalt von AT-007 (Artefakte eines einzigen Seeds); dynamische Graphen
als Quelle offener Innovation; die Identifikation von Solitonen als \emph{Quanten}-Objekte
(sie sind eine nichtlineare Verallgemeinerung von Trägern, nicht das Quantenregime); offene
Evolution in endlichen Systemen (als verzögerte Sättigung erwiesen); Symmetrie-, Attraktor-,
Topologie- und Informationsgeometrie-Wege zur Koide-Relation; Topologie-, Attraktor-,
Moduli- und Persistenz-Wege zur Eichanzahl; Stabilitäts-, Anomalie-, Darstellungs-, Defekt-
und Informations-Wege zu $N\le3$; Tired-Light-Kosmologie; und die TRM Quantum Engine.
Jedes Scheitern wurde als No-Go oder Negativresultat aufgezeichnet, und die meisten wurden
später zu den No-Go-Sätzen von Kapitel~\ref{ch:nogo} befördert.

\chapter{Warum die Struktur/Inhalt-Trennung}
\label{ch:split}

\section{Historische Motivation}

Die Struktur/Inhalt-Trennung --- \emph{Struktur ist ableitbar; Inhalt wird realisiert} ---
war kein im Voraus gewähltes Axiom. Sie war die \emph{Schlussfolgerung} eines bestimmten
Audits (QG-042, die Studie ``Parameter vs.\ Struktur''), das eine unverblümte Frage
stellte: Was lässt sich mit den Primitiven tatsächlich \emph{berechnen}, und was
\emph{ist} lediglich? Die Antwort teilte sich sauber entlang einer Linie, die zuvor
niemand explizit gezogen hatte.

Alles, was ableitbar war --- Oszillation, Phase, $U(1)$, Ladungsquantisierung,
Gravitation, Trägheit, das Äquivalenzprinzip, Teilchenstabilität, die Interpretation des
Higgs-Feldes ---, erwies sich als \emph{Form}, \emph{Symmetrie} oder \emph{Topologie}. Es
beantwortet \emph{was existiert} und \emph{warum}. Alles, was sich der Ableitung
widersetzte --- $\alpha=1/137$, $\alpha_s$, $\theta_W$, die Yukawa-Kopplungen, $\lambda$,
$\theta_{QCD}$ ---, erwies sich als \emph{dimensionslose reine Zahlen}. Sie beantworten
\emph{wie viel}. Die tiefe Trennung bildet sich exakt auf die beiden Primitive ab: $Q$
trägt die Struktur (Gesetze), und die Zufällige Aktualisierung trägt den Inhalt
(Geschichte).

Es gibt einen scharfen Grund dafür, dass die Trennung \emph{fundamental} und nicht
vorübergehend ist. Keine Kombination des dimensionbehafteten Tripels $(\ell,\tau,\hbar)$
kann eine dimensionslose Zahl erzeugen, sodass eine dimensionslose Kopplung \emph{nicht}
aus den einheitenkonventionellen Primitiven zusammengesetzt werden kann. Die Theorie
\emph{scheitert} daher nicht bloß daran, die Kopplungen abzuleiten; sie erklärt,
\emph{warum sie nicht ableitbar sind}: Sie sind historische Ergebnisse der
Aktualisierungsziehung, keine mathematischen Notwendigkeiten. Kontingenz ist intrinsisch;
kein Multiversum ist nötig. Struktur ist, was der Zufall nicht ändern kann; Inhalt ist,
was er bestimmt.

\section{Was die Trennung nicht ist}

Drei Verwechslungen müssen ausgeräumt werden, denn die feindliche Begutachtung hat alle
drei aufgeworfen. Erstens ist die Trennung \emph{keine Immunisierungsstrategie}: Sie sagt
nicht ``was immer wir nicht ableiten können, ist kontingent''. Die Grenze wird nicht im
Nachhinein verschoben; sie ist eine feste, prüfbare Teilung --- Struktur ist, was die
Primitive berechnen, Inhalt ist, was sie nicht berechnen ---, und die Theorie akzeptiert
die Konsequenzen, einschließlich der Nichtableitbarkeit ihrer eigenen Kopplungen.
Zweitens ist die Trennung \emph{kein Multiversum}: Kontingenz ist durch die Zufällige
Aktualisierung einem einzigen Universum intrinsisch, und es wird kein Ensemble von
Universen bemüht. Drittens ist die Trennung \emph{nicht anthropisch}: Die abgeleiteten
Bestandteile (räumliche $3$, $U(1)$, die Lognormalform) sind Theoreme oder Ergebnisse
über die Verteilungsform, keine Überlebensbedingungs-Selektionen, und die selektierten
Bestandteile werden als selektiert gekennzeichnet, statt hinter dem Wort ``abgeleitet''
versteckt zu werden.

Der Sinn der Trennung besteht genau darin, die Theorie \emph{in ihren Klassifikationen
falsifizierbar} zu machen. Ein mit \realunder\ gekennzeichnetes Ergebnis ist eine
stehende Einladung, die übersehene Ableitung zu finden; ein mit \drawn\ gekennzeichnetes
Ergebnis ist eine stehende Vorhersage, dass der Wert keinen tieferen Ursprung hat. Beides
sind prüfbare Behauptungen über die Grenze der Ableitbarkeit, keine rhetorischen
Immunisierungen.

\section{Beispiele}

Am leichtesten ist die Trennung anhand konkreter Fälle zu begreifen, die später jeweils
vollständig erscheinen:

\begin{itemize}
\item \textbf{Eichstruktur.} $U(1)$ (die Gruppe) ist \derived; die Farbanzahl $n=3$
(der Inhalt) ist \drawn. Die Gruppenstruktur ist aus $\mathrm{Aut}(S^1)$ berechenbar; die
Anzahl wird durch kein getestetes Prinzip festgelegt (T-09).
\item \textbf{Dimensionalität.} Die räumliche $3$ ist \derived (die
Komplexitätsmaximierung wählt $d=3{+}1$); die innere $3$ (Generationen, Farbe) ist
\emph{selektiert} --- eine abgeleitete Untergrenze $N\ge3$ trifft auf eine empirische
Obergrenze $N\le3$.
\item \textbf{Das Häufigkeitsgesetz.} Die lognormale \emph{Form} der
Häufigkeitsverteilung ist \derived (eine multiplikative Kaskade plus der zentrale
Grenzwertsatz im Log-Raum); die realisierten Werte $\mu,\sigma$ sind kontingenter Inhalt.
\item \textbf{Die Koide-Relation.} $Q=2/3$ ist real, vorhersagekräftig und RG-stabil ---
eine kanonische \realunder-Struktur ---, aber ihr Ursprung ist nicht ableitbar (T-08).
\item \textbf{Yukawa-Kopplungen.} Die Überlappoperator-\emph{Form} des Yukawa-Operators
ist abgeleitet; ihr \emph{Spektrum} (die Hierarchie $1{:}207{:}3478$) wird durch die
kontingenten Architekturformen festgelegt.
\end{itemize}

Diese Taxonomie der Nichtableitbarkeit --- \derived, \realunder, \drawn --- wird in
Kapitel~\ref{ch:taxonomy} formalisiert. Ihre Funktion ist Ehrlichkeit: Sie sagt präzise,
welche Behauptungen Theoreme und welche Klassifikationen sind, und sie ist der Grund
dafür, dass die Geschichte der feindlichen Begutachtung
(Kapitel~\ref{ch:audit-trail}) eine Geschichte der \emph{Rahmen}-Korrekturen ist und
nicht eine zurückgezogener Physik.

\chapter{Geschichte der feindlichen Begutachtung}
\label{ch:review-history}

\section{Anfängliche Kritiken}

Die Theorie wurde zwei Runden feindlicher Begutachtung unterzogen, und die Aufzeichnung
beider Runden ist Teil des Projekts. Die erste Begutachtung (des v1.0-Papiers) warf zwölf
Punkte auf; das Antwort-Audit klassifizierte sie als \emph{1 gültig}, \emph{11 teilweise
gültig} und \emph{0 ungültig}. Der dominante Fehlermodus war kein fataler Theoriebruch,
sondern eine \emph{Dokumentations}-Lücke: Das Papier ließ das dynamische System, die
formalen Primitive, das Komplexitätsfunktional, die Emergent-GR-Abgleichung und die
quantitativen Vorhersagen aus, die das Programm tatsächlich enthielt, und überzog seine
Rahmung (``geschlossen'' wo es ``disponiert'' hätte sagen sollen; ``No-Go'' wo es
``bedingtes No-Go'' hätte sagen sollen; ``Ableitung'' wo es ``ontologische
Reinterpretation'' hätte sagen sollen). Drei echte Theorielücken wurden anerkannt: das
vorläufige Eichanzahl-No-Go (T-09), der kontingente Inhalt und das Immunisierungsrisiko
der Struktur/Inhalt-Trennung.

Das Antwort-Audit identifizierte außerdem sieben Veröffentlichungsblocker, allesamt
Dokumentationskorrekturen ohne neue Physik: sechs fehlende Abschnitte (die formalen
Primitivdefinitionen, die Zusammenfassung des dynamischen Systems, das
Komplexitätsfunktional, die Emergent-GR-Ableitungszusammenfassung, die quantitativen
Vorhersagen und die Umfang-und-Grenzen-Erklärung) plus drei Formulierungskorrekturen, die
als eine einzige Rahmenkorrektur behandelt wurden. Der wichtige Punkt ist die
Klassifikation: Die Begutachtung fand \emph{einen} wirklich gültigen Einwand (die
``kein Weg offen''-Rahmung), \emph{elf} teilweise gültige Einwände, die auf ausgelassene
Dokumentation statt auf ein falsches Ergebnis zeigten, und \emph{null} vollständig
ungültige Einwände. Eine Begutachtung, die in dieser Verteilung landet, ist eine
Begutachtung der \emph{Darstellung} eines Papiers, nicht seines \emph{Kerns} --- und das
Programm behandelte sie entsprechend: Darstellung korrigieren, Kern behalten und erneut
verifizieren, dass sich nichts geändert hat.

Eine zweite Runde (Runde~2) warf zwölf ``FATAL''-Punkte auf, die in
Kapitel~\ref{ch:audit-trail} tabellarisch erfasst sind. Sie reichten von ``Primitive
als mathematische Objekte undefiniert'' bis zu ``kein kontrollierter
Kontinuumsgrenzwert.''

\section{Was sich änderte}

Die Antwort war eine Überarbeitung, die jede Rahmenübertreibung korrigierte und die sechs
fehlenden Abschnitte ergänzte (formale Primitivdefinitionen, Zusammenfassung des
dynamischen Systems, Komplexitätsfunktional, Emergent-GR-Ableitungszusammenfassung,
quantitative Vorhersagen, Umfang und Grenzen). Die Formulierungskorrekturen waren:
``geschlossen'' $\to$ ``disponiert'', ``No-Go'' $\to$ ``bedingtes No-Go'' und
``Ableitung'' $\to$ ``ontologische Reinterpretation'', wo angebracht. Keine Physik änderte
sich.

\section{Was überlebte}

Der eine substantielle Einwand, der sich nicht durch Dokumentation ausräumen ließ, war
``kein kontrollierter Kontinuumsgrenzwert.'' Er wurde auf der Schr\"odinger-Seite
\emph{vollständig} beantwortet ($L_Q\to-\nabla^2\to$ Schr\"odinger, dazu das gekrümmte
$L_W$, die Laplace--Beltrami-Näherung und die Unitarität, alles nun getestet) und auf der
Einstein-Seite \emph{teilweise} (die Standardkette funktioniert, aber die Metrik und die
BDG-Wirkung sind importiert). Der entscheidende Rest ist daher die \emph{native Metrik
$\to$ Operator-Kopplung} --- die ``G4''-Lücke: AT importiert $g_{\mu\nu}$ über den Satz
von Malament, statt es aus Q-Ereignissen zu erzeugen. Dieser Rest ist es, der das Urteil
bei \textbf{READY AS A FOUNDATION MONOGRAPH --- READY AS A RESEARCH-PROGRAM PAPER --- NOT
READY AS A DERIVATION PAPER} hält.

\chapter{Die Forschungsreise}
\label{ch:journey}

Das Ableitungsprogramm war in Forschungsprogramme gegliedert, die jeweils durch
ausführbare Experimente abgeschlossen wurden. Vier von ihnen tragen das technische
Gewicht dieser Monographie und werden hier erzählt.

\section{Das Eichprogramm}

Das Eichprogramm fragte, warum die Eichgruppe $SU(3)\times SU(2)\times U(1)$ ist. Die
Antwort ist dreigeteilt, und jeder Teil hat einen anderen Status.

\textbf{$U(1)$ ist abgeleitet.} Die Phase lebt auf dem Kreis $S^1$; ihre
Isometriegruppe \emph{ist} $U(1)$, und die Windungszahl
$\oint\nabla\theta\cdot dl=2\pi n$ liefert ganzzahlige Ladung, mit der Eichverbindung
$A_\mu=\partial_\mu\theta$ und dem Photon als der $n=0$-Phasenwelle. Dies ist das
sauberste Ergebnis des gesamten Programms: keine Wahl, sondern die Symmetrie der Phase
selbst.

\textbf{$SU(2)$ ist emergent.} Das binäre Windungspaar $\{n=\pm1\}$ ist das minimale
Dublett; die Bloch-Sphäre liefert $SO(3)=SU(2)/Z_2$; die Spinor-Doppelüberlagerung
liefert $SU(2)$. Die ``2'' ist nahezu abgeleitet, aber die Anhebung vom diskreten $Z_2$
zum kontinuierlichen $SU(2)$ ist selbst nicht abgeleitet --- sie ist die eine
zugestandene Lücke in der $SU(2)$-Geschichte.

\textbf{$SU(3)$ ist kontingent.} Das Argument $\mathrm{Aut}(C^n/S_n)\supseteq SU(n)$
leitet die Gruppen\emph{struktur} ab, aber nicht die \emph{Anzahl}, und die
8-Gluon-Algebra ist entlehnt. Vier Wegfamilien wurden für die Anzahl getestet: Topologie
allein (scheitert für $n>1$, da $\pi_1(S^1)=\mathbb Z$ unendlich ist), Attraktorraum
(scheitert, kontingent), Defektmoduli (am stärksten, aber es leitet Struktur ab, nicht
Anzahl) und Persistenz/Symmetrie (selektiert nicht). Das $1$--$2$--$3$-Muster
Rang-gleich-Windungszahl ist verlockend, aber ``möglicherweise Numerologie'': Kein
Mechanismus verbindet die drei Ränge. Das Ergebnis ist das No-Go T-09 --- kein
topologisches, Attraktor-, Moduli- oder Persistenzprinzip legt die Defektanzahl fest ---,
und das verbleibende ``warum 3'' ist die eine offene Tür der Theorie.

\subsection{Gescheiterte Wege und warum sie scheiterten}

Vier Wegfamilien scheiterten an der Defektanzahl. \emph{Topologie allein} scheitert, weil
$\pi_1(S^1)=\mathbb Z$ unendlich ist und $n=3$ nicht selektieren kann. \emph{Attraktorraum}
scheitert, weil der Landschaftsinhalt kontingent ist und daher keine Anzahl selektieren
kann. \emph{Persistenz und Symmetrie} scheitern, weil jede klassische Gruppe stabil ist,
sodass Stabilität nicht diskriminiert. Das \emph{Rang-gleich-Windungszahl}-Muster
``1--2--3'' ist verlockend, hat aber keinen verbindenden Mechanismus --- ``möglicherweise
Numerologie''. Der Defektmoduli-Weg war \emph{kein} Scheitern: Er ist der stärkste Weg,
aber er leitet die Struktur $\mathrm{Aut}(C^n/S_n)\supseteq SU(n)$ ab und hört vor der
Anzahl auf, was genau die Grenze ist, die das No-Go (T-09) markiert.

\section{Das Koide-Programm}

Die Koide-Relation $Q=\sum m/(\sum\sqrt m)^2=2/3$ ist bis auf $\sim10^{-5}$ real,
vorhersagekräftig (Vorhersage von 1981: $m_\tau=1776.97$\,MeV, 1992 bestätigt) und
RG-stabil. Das Koide-Programm fragte, woher sie kommt, und antwortete durch
\emph{Ausschöpfung}: Es probierte jeden denkbaren Weg und zeichnete auf, welche
scheiterten.

Sieben Wege wurden ausgeschöpft --- Symmetrie ($S_3$), Topologie ($S^1$), Attraktoren,
Informationsgeometrie, Gruppentheorie, der $m=3$-Abschluss und Moduli-Argumente --- mit
dem Ergebnis \emph{sechs No-Gos und einem blockierten}. Vier berechnete
Ursprungstests scheiterten alle: die demokratische $S_3$-Mischung liefert $Q=1/3$
(nicht $2/3$); es gibt keinen RG-Fixpunkt bei $2/3$; $S^1$ lokalisiert den Winkel, legt
ihn aber nicht fest; und das Informationsgeometrie-Maß $S/S_{\max}=0.5093$ ist nicht
extremal. Fünf Fehlschlüsse wurden explizit verworfen: Anthropik ($Q=2/3$ ist nicht
anthropisch), Textur-Anpassung (passt an, leitet nicht ab), Numerologie (real, aber nicht
abgeleitet), verborgene Parameter (verboten) und $45^\circ$-Umformulierungen (sie fügen
null Information hinzu).

Der Abschluss ist ein \emph{No-Go} (T-08): Der Koide-Ursprung ist eine geschlossene
Frage --- real, vorhersagekräftig, RG-stabil, \emph{nicht ableitbar} --- die kanonische
\realunder-Struktur. Das Neutrino-Koide-Analogon wurde daraufhin getestet und
\emph{falsifiziert} (T-11: $Q_{\max}=0.585$ normal / $0.500$ invertiert, beides
$<2/3$), wodurch der letzte verbleibende Unterscheider entfernt und das Programm
geschlossen wurde.

\subsection{Gescheiterte Wege und warum sie scheiterten}

Sieben Wege scheiterten, und das Muster des Scheiterns ist aufschlussreich.
\emph{Symmetrie ($S_3$)} scheitert, weil die demokratische Mischung $Q=1/3$ und ein
hierarchischer Grenzfall $Q=1$ liefert; $2/3$ ist der nicht-generische Mittelpunkt.
\emph{Topologie ($S^1$)} scheitert, weil sie die Phase lokalisiert, ohne den
$45^\circ$-Winkel festzulegen. \emph{Attraktoren} scheitern, weil es keinen RG-Fixpunkt
bei $2/3$ gibt. \emph{Informationsgeometrie} scheitert, weil $S/S_{\max}=0.5093$ nicht
extremal ist. \emph{Gruppentheorie} selektiert nicht. Der \emph{$m=3$-Abschluss}
scheitert, weil sein $\gamma\ne2/3$ ist und seine Parameter nicht auf Observablen
abgebildet sind. \emph{Moduli} ist blockiert, nicht No-Go. Die verworfenen Fehlschlüsse
--- Anthropik, Textur-Anpassung, Numerologie, verborgene Parameter,
$45^\circ$-Umformulierungen --- scheiterten, weil jeder entweder anpasst, ohne
abzuleiten, oder umformuliert, ohne Information hinzuzufügen.

\section{Das Kontinuumsprogramm}

Das Kontinuumsprogramm fragte, ob das diskrete dynamische System einen kontrollierten
Kontinuumsgrenzwert besitzt, und es zerfiel in zwei Hälften mit ungleichen Ergebnissen.

Auf der Schr\"odinger-Seite ist die Antwort ein eindeutiges \emph{Ja}. Der
Ketten-Laplace-Operator $L_Q$ konvergiert gegen den flachen Laplace-Operator $-\nabla^2$
mit einem exakten Spektrum in geschlossener Form und Konvergenz zweiter Ordnung, und der
gewichtete Laplace-Operator $L_W$ liefert eine gekrümmte Verallgemeinerung und eine
Laplace--Beltrami-Näherung. Diese Hälfte ist vollständig getestet.

Auf der Einstein-Seite ist die Antwort zurückhaltender. Der BDG-Kausalmengen-Operator
konvergiert gegen den d'Alembert-Operator $\square$ mit $O(h^2)$-Kontrolle, aber $L_Q$
und BDG sind \emph{in der Signatur disjunkt} --- der eine elliptisch und positiv, der
andere hyperbolisch und indefinit ---, sodass keine native Brücke zwischen ihnen
geschlossen wurde. Ein eigenes Drei-Test-Programm machte die Inkompatibilität ausführbar:
$L_Q$ ist positiv semidefinit, BDG ist indefinit, und der Graphen-Laplace-Operator wird
für Lorentz-Signatur explizit verworfen. Dies --- zwei disjunkte Ketten auf demselben
Substrat --- ist die wichtigste strukturelle Tatsache des Programms und der präzise Grund
dafür, dass die Einstein-Wiederherstellung logisch statt mathematisch ist.

\subsection{Gescheiterte Wege und warum sie scheiterten}

Der eine Weg, der scheiterte --- und der Grund, warum er scheiterte ---, ist die naive
Hoffnung, ein einzelner Operator könne sowohl der Schr\"odinger- als auch der
Einstein-Seite dienen. Er scheiterte, weil die beiden Kontinuumsgrenzwerte inkompatible
Signaturen haben: $L_Q\to-\nabla^2$ ist elliptisch und positiv semidefinit, während
$\mathrm{BDG}\to\square$ hyperbolisch und indefinit ist. Der ausführbare Test machte das
Scheitern präzise: Der Graphen-Laplace-Operator wird für Lorentz-Signatur explizit
verworfen. Das Programm erzeugte daher keine fehlende ``Brücke''; es erzeugte ein
\emph{Kein-Brücke-Ergebnis} --- zwei disjunkte Ketten auf demselben Substrat ---, das die
wichtigste strukturelle Tatsache des Programms ist.

\section{Das Metrikprogramm}

Das Metrikprogramm fragte, wie eine Metrik $g_{\mu\nu}$ überhaupt entstehen kann. Die
Antwort ist eine viergliedrige Kette: Q-Ereignisse $\to$ Kausalordnung (nativ) $\to$
konforme Klasse (bewiesen, über den importierten Satz von Malament) $\to$ konformer
Faktor (nativ, aus dem Zählmaß $f=\rho^{2/d}$) $\to$ die Metrik. Der
Metrik\emph{ursprung} ist daher geschlossen: Das einzige importierte Glied ist ein
bewiesener Satz, keine Theorielücke. Die Verifikation ist direkt --- die Kausalordnung
erfüllt die Metrik-Axiome, die Nullstruktur ist unter konformer Umskalierung invariant,
aber nicht unter nicht-konformer Umskalierung, und das Zählmaß reproduziert den konformen
Faktor im Rundlauf.

Was offen bleibt, ist die Metrik\emph{dynamik} --- die native Metrik $\to$
Operator-Kopplung (G4). Den Ursprung zu schließen sagt uns, \emph{woher} die Metrik
kommt; es sagt dem Operator nicht, wie er sie \emph{verwenden} soll. Dieser Rest ist der
Grund, warum die Einstein-Wiederherstellung logisch statt mathematisch bleibt, und er ist
der einzige wirklich offene technische Punkt der Monographie.

\subsection{Gescheiterte Wege und warum sie scheiterten}

Der Weg, der scheiterte, ist die \emph{native} Metrikdynamik. AT erzeugt $g_{\mu\nu}$
nicht aus Q-Ereignissen; der metrikabhängige Operator $\Delta_g/\Box_g$ fehlt, und der
gewichtete Laplace-Operator $L_W$ liefert nur eine Näherung, deren definierende Metrik
extern ist. Der Versuch, die Metrik\emph{dynamik} nativ zu schließen, scheiterte daher,
und der ehrliche Endpunkt des Programms ist die Trennung, die es festhält: Der
Metrik\emph{ursprung} ist geschlossen (Q-Ereignisse $\to$ Kausalordnung $\to$ konforme
Klasse $\to$ konformer Faktor $\to$ Metrik, wobei der Satz von Malament als bewiesenes
Glied importiert ist), während die Metrik\emph{dynamik} importiert bleibt. Diese Trennung
ist die ``G4''-Lücke.

\section{Was die vier Programme gemeinsam errichten}

Zusammen gelesen zeichnen die vier Programme einen Bogen. Das Eich- und das
Koide-Programm lokalisieren die \emph{Grenze} zwischen dem Ableitbaren und dem
Kontingenten; das Kontinuums- und das Metrikprogramm lokalisieren die \emph{Grenze}
zwischen dem Kontrollierten und dem Importierten. Beide Grenzen sind ehrlich: Die Theorie
behauptet, Struktur abzuleiten und den Schr\"odinger-Grenzwert zu kontrollieren, und sie
verzichtet ausdrücklich darauf, die Eichanzahl, den Koide-Ursprung oder die native
Einstein-Brücke zu beanspruchen. Diese Teilung --- ableitbare Struktur, kontrollierter
Schr\"odinger-Grenzwert, klassifizierter Inhalt, importierte Metrik --- ist die Gestalt
der fertigen Theorie.

\chapter{Gewonnene Lehren}
\label{ch:lessons}

\section{Falsche Annahmen}

\begin{itemize}
\item \textbf{Der statische Graph war die verborgene Annahme Nummer eins.} Das erste
Audit alternativer Grundlagen (X001) fand elf verborgene Annahmen, mit dem statischen
Graphen an erster Stelle der Liste. Als er getestet wurde (X002--X004), erzeugten
dynamische Graphen keine neuen Phänomene und keine offene Innovation.
\item \textbf{Nichtlinearität zerbricht das Eigenmoden-Bild.} X005 zeigte, dass die
meisten frühen ``Quanten''-Ergebnisse Artefakte der linearen Algebra waren:
Nichtlinearität zerstört Eigenmoden und erzeugt stattdessen Solitonen. Q, Fitness und
Evolution überlebten; die naive Quanten-Lesart nicht.
\item \textbf{Offene Evolution in endlichen Systemen ist eine Illusion.} X025 schien
offenes Artenwachstum zu zeigen, aber X026--X027 bewiesen, dass es verzögerte Sättigung
war: Ein Schubfachargument begrenzt die Artenzahl durch die Zahl der Knoten, sodass ein
endliches System stets sättigt. Evolution der Stufe 6 (offen) erfordert einen
unendlichen Zustandsraum.
\item \textbf{Scheinbare Vielfalt kann ein einziger Keim sein.} Die
AT-007-Resonanz-``Familien'' kollabierten in einer Reproduzierbarkeitsstudie mit 400
Simulationen auf einen universellen Modus (F4); der Rest waren Artefakte.
\end{itemize}

\section{Verworfene Ableitungen}

Das Programm versuchte wiederholt --- und verwarf --- Kandidatenableitungen, statt sich
mit Postdiktion zufriedenzugeben: Symmetrie-, Attraktor-, Topologie- und
Informationsgeometrie-Wege zu Koide (T-08); Topologie-, Attraktor-, Moduli- und
Persistenz-Wege zur Eichanzahl (T-09); Stabilitäts-, Anomalie-, Darstellungs-, Defekt-
und Informations-Wege zu $N\le3$ (T-10); fünf getrennte Versuche, $\alpha=1/137$
abzuleiten; die Tired-Light-Kosmologie; die TRM Quantum Engine; und die großen
vereinheitlichten Gruppen $SU(5)/SO(10)$ (Protonenzerfall unbeobachtet) und $E_6/E_8$
(exotische Teilchen unbeobachtet). In jedem Fall wurde die Verwerfung aufgezeichnet, und
wo das Scheitern systematisch war, wurde daraus ein No-Go-Satz.

\section{No-Go-Entdeckungen}

Die No-Go-Sätze (Kapitel~\ref{ch:nogo}) sind das markanteste wissenschaftliche Produkt
des Programms. T-08 (Koide nicht ableitbar), T-09 (Eichanzahl nicht festgelegt), T-10
($N\le3$ nicht festgelegt), T-11 (Neutrino-Koide falsifiziert) und T-12 (gemeinsame
Kaskade untestbar) definieren gemeinsam die Grenze zwischen dem, was die Theorie
berechnen kann, und dem, was sie klassifizieren muss. Die Meta-Lehre ist die
Struktur/Inhalt-Trennung selbst: Die No-Gos sind keine Einfallsfehler, sondern die
erwartete Signatur einer Theorie, in der Inhalt konstruktionsbedingt kontingent ist.

\section{Meta-Lehren}

Jenseits der spezifischen falschen Annahmen und verworfenen Ableitungen ergab das Programm
eine kleine Zahl dauerhafter Lehren, die alles danach prägten.

\textbf{Die Ableitbarkeitsgrenze ist selbst ein Ergebnis.} Die Entdeckung, dass Struktur
ableitbar ist, während dimensionsloser Inhalt es nicht ist (QG-042), ist keine
Beschränkung, in die das Programm hineinstolperte; sie ist ein positiver Befund und
lizenziert die Taxonomie. Eine Theorie, die \emph{erklärt, warum} ihre Kopplungen nicht
ableitbar sind, erhebt einen stärkeren Anspruch als eine Theorie, die sie bloß nicht
ableiten kann.

\textbf{Negative Ergebnisse sind Produkte erster Klasse.} Die No-Go-Sätze (T-08--T-12)
sind der markanteste wissenschaftliche Ertrag des Programms, nicht seine Peinlichkeiten.
Ein Weg, der als geschlossen erwiesen wird, ist Information, und ihn als Satz
festzuhalten (statt ihn still fallenzulassen) macht die Theorie prüfbar.

\textbf{Ehrlichkeit in der Rahmung ist ein wissenschaftliches Gut.} Die Geschichte der
feindlichen Begutachtung ist fast vollständig eine Geschichte von Rahmenkorrekturen, nicht
zurückgezogener Physik. Die Lehre ist, dass ``geschlossen'' vs.\ ``disponiert'',
``No-Go'' vs.\ ``bedingtes No-Go'' und ``Ableitung'' vs.\ ``ontologische
Reinterpretation'' keine kosmetischen Unterscheidungen sind; sie sind der Unterschied
zwischen einer Überbeanspruchung und einer verteidigungsfähigen Aussage.

\textbf{Ausführbare Verifikation macht Ehrlichkeit überprüfbar.} Jedes ``BESTANDEN'' in
dieser Monographie ist das Ergebnis eines Tests, den ein feindlicher Gutachter ausführen
kann.
Diese Disziplin verwandelte die Einwände der feindlichen Begutachtung in einen
Auditpfad, der Punkt für Punkt beantwortet werden konnte, statt endlos diskutiert zu
werden.

\chapter{Evolution des Rahmens}
\label{ch:evolution}

Der Rahmen erreichte seine endgültige Form nicht in einem Schritt; er wurde wiederholt
abgebaut und neu aufgebaut. Dieses Kapitel zeichnet nach, was zu Beginn angenommen, was
entfernt, was beibehalten wurde und welche zwei methodischen Verpflichtungen entstanden
--- die Struktur/Inhalt-Trennung und die No-Go-Methode.

\section{Früheste Annahmen}

Das Programm begann mit einem reichen und größtenteils falschen Satz von Annahmen. Materie
wurde als dynamisch stabilisierte Wellenstruktur in einem temporalen Feld angenommen;
Synchronisation und Resonanz wurden als fundamentaler als Teilchen betrachtet; und der Weg
von den Oszillatoren zur Kosmologie wurde als eine einzige durchgehende Kette angenommen.
Implizit nahmen die frühen Experimente außerdem an, dass ein statischer Graph das richtige
Substrat sei, dass Zufallsmatrizen für Emergenz ausreichen könnten und dass jede
auftretende kohärente Struktur die Quanten-Lesart des Systems tragen würde.

\section{Entfernte Annahmen}

Die meisten frühesten Annahmen wurden experimentell entfernt. Zufallsmatrizen erzeugten
keine Emergenz (AT-002); statische Topologie allein war unzureichend (AT-003); die
scheinbare Vielfalt der Resonanzfamilien kollabierte auf einen einzigen reproduzierbaren
Modus (AT-008); dynamische Graphen erzeugten keine offene Innovation (X002--X004); und
Nichtlinearität zeigte, dass die frühen ``Quanten''-Ergebnisse Artefakte der linearen
Algebra waren (X005). Die spezifischen Annahmen des Vorgänger-TRM-Rahmens ---
Tired-Light-Kosmologie und eine nicht-unitäre Quantenmaschine --- wurden im
Legacy-Audit rundheraus verworfen. Was nach diesem Abstreifen blieb, war minimal:
Individuation ($Q$), Reversibilität, Selbstkonsistenz und echter Zufall.

\section{Beibehaltene Annahmen}

Zwei Paare überlebten jede feindliche Reduktion. \emph{Reversibilität} (erhaltene Norm)
und \emph{Selbstkonsistenz} (Fixpunkte) erwiesen sich als unabhängig und minimal
hinreichend: Ihr Schnitt ist die unitäre Quantenmechanik (X011--X015). \emph{$Q$}
(Individuation) und \emph{echter Zufall} erwiesen sich als irreduzibel --- unter ihnen
existiert keine tiefere Invariante (X035, X039). Alles andere in der Theorie --- Zeit,
Raum, $3{+}1$, die Schr\"odinger-Gleichung, die Bornsche Regel, Messung, Gravitation als
Kausalstruktur, Teilchen als Defekte und $U(1)$ --- ist aus diesen abgeleitet oder
emergent, plus der einzelnen kontingenten Zahl $M^2$.

\section{Entstehung der Struktur/Inhalt-Trennung}

Die Struktur/Inhalt-Trennung wurde nicht auferlegt; sie wurde \emph{entdeckt}. Das
Parameter-vs.-Struktur-Audit (QG-042) fand, dass alles aus den Primitiven Ableitbare
\emph{Form}, \emph{Symmetrie} oder \emph{Topologie} war, während alles, was sich der
Ableitung widersetzte, eine \emph{dimensionslose reine Zahl} war. Da keine Kombination des
dimensionbehafteten Tripels $(\ell,\tau,\hbar)$ eine dimensionslose Zahl erzeugen kann,
\emph{können} die Kopplungen nicht aus den Primitiven zusammengesetzt werden. Das
verwandelte eine vorübergehende Beschränkung in ein Prinzip: Struktur ist ableitbar,
Inhalt wird realisiert, und die Taxonomie (\derived, \realunder, \drawn) ist die ehrliche
Buchführung dieser Grenze.

\section{Entstehung der No-Go-Methodik}

Die No-Go-Methode entstand aus derselben Disziplin. Wenn eine Ableitung scheiterte ---
Koide, die Eichanzahl, $N\le3$ ---, verwarf das Programm das Scheitern nicht; es zählte
jeden denkbaren Weg auf, testete jeden gegen die Keine-neuen-Primitive-Beschränkung und
hielt das Ergebnis als bedingten No-Go-Satz fest (T-08--T-12). Was als Methode begann,
einzelne Fragen zu schließen, wurde zu einer allgemeinen Methode: \emph{Weg-Ausschöpfung}
plus \emph{feindliche Reduktion}, die ``wir konnten es nicht ableiten'' in die stärkere,
überprüfbare Aussage ``kein getesteter Weg leitet es ab'' verwandelt.

\chapterend
{Der Rahmen wurde von einem reichen Oszillatorbild auf zwei Primitive plus eine Zahl
abgestreift; die Struktur/Inhalt-Trennung und die No-Go-Methode entstanden als seine zwei
organisierenden Verpflichtungen.}
{Ob eine native Metrik-Operator-Kopplung (G4), eine eindeutige diskriminierende
Vorhersage oder eine native BDG-Wirkung ableitbar ist --- die drei
Ableitungspapier-Blocker.}
{Historisch/redaktionell; kein ausführbarer Test greift, aber jede hier zitierte
Entfernung und Beibehaltung ist im Entwicklungsregister (Anhang) festgehalten.}

\chapter{Die Methode ausführbarer Audits}
\label{ch:method}

Das markanteste Merkmal dieses Programms ist kein bestimmtes Ergebnis, sondern eine
\emph{Methode}: Jede substantielle Behauptung wird ausführbar gemacht, und jedes
``-Audit'' in der Aufzeichnung ist ein Stück dieser Methode. Dieses Kapitel erklärt die
Maschinerie, denn die durch die technischen Teile verstreuten Konfidenzaussagen sind ohne
sie unverständlich.

\section{Was ein Audit ist}

Ein Audit ist eine strukturierte, reproduzierbare Untersuchung mit einer einzigen Frage
und einer Disposition am Ende. Es ist als Analysator (ein kleines Programm) plus einem
oder mehreren Tests implementiert und schreibt einen Bericht. Die Frage wird zuerst
genannt, die Wege werden aufgezählt, und jeder Weg wird gegen einen festen Satz von
Beschränkungen getestet --- typischerweise die \emph{Keine-neuen-Primitive}-Beschränkung
und eine Verwerfung von Anthropik, Numerologie, verborgenen Dimensionen und
Post-Selektion, sofern sie nicht ausdrücklich benannt sind. Die Disposition ist eines aus
einem kleinen Vokabular: \derived, \realunder, \drawn, \textbf{No-Go}, \textbf{blockiert},
\textbf{kontingent} oder \textbf{selektiert}. Weil der Bericht in eine Datei geschrieben
wird und der Test auf den berechneten Zahlen prüft, kann ein späterer Gutachter das Audit
erneut ausführen und dieselbe Antwort erhalten.

\section{Weg-Ausschöpfung und die No-Go-Methode}

Die No-Go-Sätze werden durch \emph{Weg-Ausschöpfung} erzeugt: Zähle jeden
Ableitungsweg auf, den das Programm für ein Ziel ersinnen kann, teste jeden und zeichne
auf, welche scheitern. Ein No-Go ist die Aussage ``innerhalb des getesteten Wegraums und
der Keine-neuen-Primitive-Beschränkung erreicht kein Weg das Ziel.'' Es ist daher
\emph{bedingt} --- es ist kein Beweis logischer Unmöglichkeit (außer wo eine Berechnung,
wie T-11, die Frage direkt schließt), und es ist nur so stark wie der aufgezählte
Wegraum. Der Koide-Abschluss ist das Paradigma: sieben Wege, sechs No-Gos, einer
blockiert, null offen, und die Schlussfolgerung ist eine \emph{geschlossene Frage}, keine
Ableitung.

\section{Konfidenzränge}

Die Konfidenzzahlen sind \emph{unkalibrierte ordinale Ränge}. Sie werden vom Audit
vergeben, um den eigenen Überzeugungsgrad des Programms in eine Disposition zu ordnen,
und tragen absichtlich keine statistische Bedeutung. Ein Satz mit vollständiger
Weg-Ausschöpfung hinter sich (wie $U(1)$, Konfidenz 0,95) steht über einer emergenten
Struktur mit zugestandener Lücke ($SU(2)$, 0,70), die über einem vorläufigen No-Go
($SU(3)$-Struktur, 0,10) steht. Die Zahl ist eine Zusammenfassung davon, \emph{wie viel
des Wegraums geschlossen wurde}, keine A-posteriori-Wahrscheinlichkeit.

\section{Feindliche Reduktion als Methode}

Das Programm behandelt die feindliche Begutachtung nicht als äußeres Ereignis, sondern
als ständige interne Methode: Jedes Ergebnis wird periodisch einer \emph{feindlichen
Reduktion} unterzogen, die versucht, es wegzuableiten, wegzuerklären oder zu zeigen, dass
es eine Umformulierung ist. Die überlebenden Ergebnisse sind genau diejenigen, die die
feindliche Reduktion nicht entfernen kann. Deshalb ist die Geschichte der feindlichen
Begutachtung (Kapitel~\ref{ch:review-history}) so arm an zurückgezogener Physik: Die
Rahmung wurde korrigiert, aber die zugrunde liegenden Ergebnisse hatten die feindliche
Reduktion intern bereits durchlaufen, bevor der externe Gutachter sie sah.

%=============================================================================
\part{Grundlagen}
%=============================================================================

\chapter{Die Primitive}
\label{ch:primitives}

Dieses Kapitel formuliert die Theorie. Alles, was in den technischen Teilen folgt, ist
eine Entfaltung dieser Definitionen: Die Primitive sind die gesamte Eingabe, und die Last
der Monographie besteht darin zu zeigen, wie viel beobachtbare Struktur aus ihnen folgt,
und ehrlich zu sagen, wo sie aufhört.

\section{Formulierung der Theorie}

THE Q-MODEL beruht auf einer einzigen organisierenden Behauptung, der
\textbf{Struktur/Inhalt-Trennung}:

\begin{center}
\emph{Struktur ist ableitbar; Inhalt wird realisiert.}
\end{center}

Die \emph{Form} beobachtbarer Struktur (Topologie, Symmetrie, Verteilungsform, untere
Schranken) ist aus einer minimalen Primitivmenge beweisbar. Der \emph{Inhalt} (spezifische
Massen, Kopplungen, Multiplizitäten, Winkel) ist eine Ziehung aus einem kontingenten
Ensemble und wird daher klassifiziert statt abgeleitet.

\section{Die vier Primitive}

\begin{definition}[Primitivmenge]
Die Theorie lässt genau vier Primitive zu, und keine weiteren:
\begin{enumerate}
\item \textbf{$Q$} — das Individuationsprinzip (Ontologie): der irreduzible Akt, durch
den ein diskretes Ereignis aus dem Nichts individuiert wird.
\item \textbf{Zufällige Aktualisierung} — echter ontologischer Zufall: Gegeben $Q$, sind
die realisierten Ereignisorte innerhalb der Kausalstruktur zufällig. (Annahme
\textbf{A-03}.)
\item \textbf{$(\ell,\tau,\hbar)$} — das irreduzible physikalische Tripel: eine
Raumzeit-Skala $\ell$, eine Uhr $\tau$ und eine Wirkung $\hbar$
(Einheitenkonventionen).
\item \textbf{$M^2$} — der einzige kontinuierliche Nichtlinearitätsparameter (das
Nichtlinearitätsregime der Dynamik).
\end{enumerate}
\end{definition}

Kein Eichgruppen-, Multiplizitäts- oder Verborgene-Dimensionen-Primitiv ist erlaubt. Die
Primitive haben zwei Schichten, die sich bei $Q$ treffen: eine \emph{Ontologie}-Schicht
(die die Ableitung von Struktur trägt) und eine \emph{dynamische} Schicht (die das
dynamische System von Kapitel~\ref{ch:dynamics} trägt).

\begin{remark}
$Q$ individuiert; die Zufällige Aktualisierung realisiert; $(\ell,\tau,\hbar)$ legt
Einheiten fest; $M^2$ setzt den einzigen kontinuierlichen Parameter, den die
Ableitungshierarchie auf $M^2\approx5$ festnagelt (Kapitel~\ref{ch:complexity}).
\end{remark}

\section{Die Quantenpostulate (dynamische Schicht)}

\begin{axiom}[Q existiert]
Topologische Ladungsquanten besitzen Position $x_i$ und Phase $\theta_i$ und
wechselwirken paarweise mit Kopplung $J_{ij}=f(|x_i-x_j|)$. (Angenommen, irreduzibel.)
\end{axiom}

\begin{axiom}[Reversible Dynamik]
Die Zustandsnorm $\|\psi\|^2$ ist erhalten; äquivalent dazu ist die Evolution unitär.
\end{axiom}

\begin{axiom}[Bornsche Regel]
$P=|\langle\phi|\psi\rangle|^2$, eindeutig selektiert durch Additivität (Satz von
Gleason, 1957).
\end{axiom}

\begin{axiom}[Messung]
Eine Messung liefert ein einziges bestimmtes Ergebnis (das Kollaps-Axiom).
\end{axiom}

Die Postulate 1--2 leiten den Hilbertraum und die Schr\"odinger-Gleichung ab
(Kapitel~\ref{ch:dynamics}); die Postulate 3--4 schließen die Interpretation.

%=============================================================================

\chapterend
{Die vier Primitive ($Q$, Zufällige Aktualisierung, $(\ell,\tau,\hbar)$, $M^2$) und die
Struktur/Inhalt-Trennung sind formuliert; kein weiteres Primitiv ist erlaubt.}
{Das Maß und das Wirkungsfunktional für $Q$ bleiben Postulate, keine abgeleiteten
Objekte.}
{Auf Axiome greift kein ausführbarer Test; die Formalisierungs-Audits im nächsten Kapitel
halten den Objekt-/Zustands-/Operations-Status fest.}

\chapter{Formalisierung von $Q$ und der Zufälligen Aktualisierung}
\label{ch:formalization}

Dieses Kapitel hält den formalen Status der Primitive fest, wie er durch die
Formalisierungs-Audits (ausführbar; siehe Teil~V) festgestellt wurde. Sein Zweck ist es,
präzise zu sagen, wie viel jedes Primitivs bereits ein mathematisches Objekt ist und wie
viel ein Postulat bleibt --- die ehrliche Grenze, auf die die feindliche Begutachtung
gedrängt hat.

\section{Status von $Q$}

\begin{center}
\begin{tabular}{lll}
\toprule
Komponente & Status & Inhalt \\
\midrule
Objekt & \textbf{Formalisiert} & ein diskretes Ereignis (individuierte Entität) \\
Zustandsraum & \textbf{Formalisiert} & Positionen $x_i$ und Phasen $\theta_i$ \\
Operationen & \textbf{Formalisiert} & Individuation, Kopplung $J_{ij}$ \\
Konfigurationsraum & \textbf{Teilweise} & Graph von Ereignissen, kein nativer Kontinuumsgrenzwert \\
Dynamik & \textbf{Teilweise} & $L_Q$-Generator (Postulat 1) \\
Symmetrien & \textbf{Teilweise} & $S^1$-Phase, $Z_2$, Permutation $S_n$ \\
Kontinuumsgrenzwert & \textbf{Teilweise} & für $L_Q$ verifiziert (Teil~II) \\
Maß & \textbf{Fehlt} & kein Maßraum allein aus den Primitiven \\
\bottomrule
\end{tabular}
\end{center}

\textbf{Klassifikation:} $Q$ ist \emph{teilweise formalisiert} — Objekt, Zustandsraum und
Operationen existieren; das Maß und das Wirkungsfunktional bleiben Postulate.

\section{Status der Zufälligen Aktualisierung}

Die Zufällige Aktualisierung ist \textbf{A-03}, eine Annahme, kein abgeleitetes Objekt.
Das Formalisierungs-Audit fand:

\begin{center}
\begin{tabular}{lll}
\toprule
Komponente & Status & Inhalt \\
\midrule
Zufallsvariable & \textbf{Formalisiert} & die Ziehung der Ereignisorte \\
Generator & \textbf{Formalisiert} & ontologischer Zufall \\
$(\Omega,\mathcal F,P)$ & \textbf{Teilweise} & Ensemble nicht spezifiziert \\
Ensemblemaß & \textbf{Teilweise} & kein explizites Maß \\
\bottomrule
\end{tabular}
\end{center}

\textbf{Klassifikation:} Die Zufallsvariable und der Generator sind formalisiert; der
Wahrscheinlichkeitsraum $(\Omega,\mathcal F,P)$ und das Ensemblemaß sind nur teilweise
formalisiert.

\begin{remark}
Der Vorwurf des feindlichen Gutachters, ``die Zufällige Aktualisierung sei Verbalismus'',
ist daher \emph{teilweise gültig}: Sie ist ein \emph{zugestandenes} Primitiv (A-03), kein
abgeleitetes Objekt, und das Papier stellt dies ausdrücklich fest.
\end{remark}

%=============================================================================

\chapter{Das dynamische System}
\label{ch:dynamics}

Dieses Kapitel macht aus den Primitiven eine Dynamik. Der folgenreichste Zug ist die
Identifikation des Graphen-Laplace-Operators $L_Q$ als Generator, denn seine Eigenbasis
\emph{ist} der Hilbertraum der Theorie, und sein Kontinuumsgrenzwert trägt das gesamte
Schr\"odinger-Programm.

\section{Der Graphen-Laplace-Operator $L_Q$}

\begin{definition}[Graphen-Laplace-Operator]
Für den $Q$-Wechselwirkungsgraphen mit Adjazenzmatrix $A$ und Gradmatrix $D$ gilt
\begin{equation}
L_Q = D - A .
\end{equation}

\eqnotes
{Jede Zeile summiert sich zu null, sodass der uniforme Zustand die Nullmode ist und die
Eigenvektoren die stehenden Moden des Graphen sind.}
{Er ist der diskrete Diffusions-/Hüpfoperator der Theorie.}
{Seine Eigenbasis \emph{ist} der Hilbertraum, und sein Kontinuumsgrenzwert ist der
Laplace-Operator, der das gesamte Schr\"odinger-Programm trägt.}
{Hier ungerichtet und ungewichtet; der Lorentz-Fall (BDG) und der gekrümmte Fall ($L_W$)
sind separate Operatoren.}
$L_Q$ ist reell, symmetrisch, positiv semidefinit und hat Null-Zeilensummen. Seine
Eigenvektoren bilden eine Orthonormalbasis — den \textbf{Hilbertraum} der Theorie.
\end{definition}

\section{Die Tight-Binding-Identität (AT-142)}

\begin{proposition}[Tight-Binding-Identität]
Auf einer 1D-Kette gilt
\begin{equation}
(L_Q)_{ij} = 2\delta_{ij} - \delta_{i,j+1} - \delta_{i,j-1},
\end{equation}
was \emph{identisch} mit dem Tight-Binding-Hamiltonian mit On-Site-Energie
$\varepsilon=2t$ ist: $H = t L_Q$. Dies ist eine mathematische \textbf{Identität}, keine
Analogie.
\end{proposition}

\theoremnotes
{Auf einer Kette ist der Laplace-Operator exakt der Zweite-Differenzen-Operator, sodass
``Graphendiffusion'' und ``Tight-Binding-Hüpfen'' dasselbe Objekt sind, zweimal
hingeschrieben.}
{Die Identität lizenziert die Übertragung von Festkörpertechniken --- Spektren,
Greensche Funktionen, Bandstruktur --- ohne jede Übersetzung auf die Ontologie-Schicht.}
{Sie gilt in einer Dimension und für den nackten (ungewichteten) Laplace-Operator; die
gekrümmten und kausalmengenartigen Verallgemeinerungen sind separate Objekte (Abschnitte
\ref{ch:laplace-beltrami} und~\ref{ch:bridge}).}

\section{Schr\"odinger aus Reversibilität (AT-149--151)}

\begin{theorem}[Schr\"odinger-Gleichung aus Reversibilität]
Betrachte lineare Evolution $\partial_t\psi = M\psi$ mit erhaltener Norm
$\partial_t\|\psi\|^2 = \psi^\dagger(M+M^\dagger)\psi = 0$. Dann gilt
\begin{equation}
M^\dagger = -M .
\end{equation}
Für reelles $M$ ist dies $M^T=-M$. Die einfachste antisymmetrische $2\times2$-Matrix $J$
erfüllt $J^2=-I$ (sodass $J$ als imaginäre Einheit wirkt), und
\begin{equation}
M = J\otimes L_Q \quad\Longrightarrow\quad i\partial_t \psi = L_Q\,\psi,
\end{equation}

\eqnotes
{Normerhaltung zwingt den Generator, antihermitesch zu sein; der Block $J$ mit $J^2=-I$
liefert die imaginäre Einheit.}
{Es ist die unitäre Dynamik der Theorie: die Schr\"odinger-Gleichung.}
{Sie bedeutet, dass die Quantenmechanik aus einer Erhaltungsanforderung \emph{erzeugt}
wird, nicht postuliert.}
{Linear und Einteilchen; die imaginäre Struktur wird über $J$ importiert.}
mit unitärer Evolution
\begin{equation}
\psi(t) = e^{-iL_Qt}\,\psi(0).
\end{equation}
\end{theorem}

\theoremnotes
{Die Erhaltung der Norm zwingt den Generator, antihermitesch zu sein; das einfachste
antihermitesche Objekt auf dem Graphen ist $J\otimes L_Q$, und $J^2=-I$ ist die imaginäre
Einheit in Verkleidung. Das ``$i$'' der Quantenmechanik ist also nicht postuliert,
sondern aus einer $2\times2$-antisymmetrischen Matrix gebaut.}
{Dies ist der dynamische Kern der Theorie: Die Schr\"odinger-Gleichung ergibt sich aus
einer Erhaltungsanforderung plus dem Graphen-Laplace-Operator, statt als Axiom behauptet
zu werden.}
{Sie ist linear und Einteilchen. Die imaginäre Struktur wird über den $2\times2$-Block $J$
importiert, und das Argument rekonstruiert für sich genommen nicht den feldtheoretischen,
wechselwirkenden Inhalt der Quantenmechanik.}

\historynotes
{Die Schr\"odinger-Gleichung wird normalerweise \emph{postuliert}; das
Reduktionsprogramm fragte, ob ihr ``$i$'' und ihre unitäre Evolution stattdessen aus
einer Erhaltungsanforderung abgeleitet werden könnten.}
{Wenn der dynamische Kern abgeleitet statt angenommen werden könnte, würde die Theorie
Quantenmechanik \emph{erzeugen}, statt sie zu importieren --- ein grundlegender, kein
kosmetischer Gewinn.}
{Der Tight-Binding-/Hilbertraum-Weg verwendete $L_Q$ als Generator, aber die spezifische
Reduktion der imaginären Einheit auf einen antisymmetrischen $2\times2$-Block $J$ mit
$J^2=-I$ war der Schritt, der noch nicht ausgeschöpft war.}
{Akzeptiert als Satz (AT-149--151), weil er exakt ist und weil sein Kontinuumsgrenzwert
dann unabhängig durch \texttt{GraphLaplacianContinuumTests} verifiziert wird.}

\section{Observablen aus dem Spektrum (AT-145)}

\begin{center}
\begin{tabular}{ll}
\toprule
Observable & Definition \\
\midrule
effektive Masse & $m_{\rm eff} = 1/\lambda_1$ \\
Energie & $E = \mathrm{tr}(L)$ \\
Lücke & $\Delta = \lambda_2-\lambda_1$ \\
Korrelationslänge & $\xi = 1/\sqrt{\lambda_1}$ \\
Diffusionskonstante & $D = \lambda_1$ \\
Kanalkapazität & $C = \log_2 N$ \\
\bottomrule
\end{tabular}
\end{center}

\begin{remark}
Dies sind \emph{dimensionale Zuweisungen}, keine dynamischen Vorhersagen; der feindliche
Gutachter merkt zu Recht an, dass dies eine Umfangsgrenze ist.
\end{remark}

\section{Die Ontologie $\to$ Physik-Brücke}

Die Brücke hat zwei Pfeiler, beide nun getestet:
\begin{enumerate}
\item $L_Q$-Eigenvektorbasis $=$ Hilbertraum (AT-149); und
\item der Kausalmengen-Kontinuumsgrenzwert liefert die Metrik (QG-001, XC006--012).
\end{enumerate}

%=============================================================================
\part{Das Kontinuumsgrenzwert-Programm (verifiziert)}
%=============================================================================

Dieser Teil hält die beiden Kontinuumsketten der Theorie und ihre (partielle) Verbindung
fest. Die eine Kette ist \emph{elliptisch}: Der Graphen-Laplace-Operator $L_Q$ konvergiert
gegen den flachen Laplace-Operator $-\nabla^2$, und seine gewichtete Verallgemeinerung
liefert eine gekrümmte Schr\"odinger-Gleichung und eine Laplace--Beltrami-Näherung. Die
andere ist \emph{Lorentz-artig}: Der Kausalmengen-BDG-Operator konvergiert gegen den
d'Alembert-Operator $\square$. Die beiden Ketten sind in der Signatur disjunkt, und die
zentrale Tatsache des Teils ist, dass zwischen ihnen bislang keine native Brücke
geschlossen wurde. Die Schr\"odinger-Seite ist kontrolliert und getestet; die
Einstein-Seite ist partiell.

\chapterend
{$L_Q=D-A$ ist der Generator und seine Eigenbasis der Hilbertraum; die
Schr\"odinger-Gleichung $i\partial_t\psi=L_Q\psi$ folgt aus der Normerhaltung.}
{Die spektrumsabgeleiteten Observablen sind dimensionale Zuweisungen, keine dynamischen
Vorhersagen.}
{Mathematisch (AT-142 Tight-Binding-Identität; AT-149--151 Reversibilitätssatz); die
Kontinuumsgrenzwert-Tests im nächsten Teil verifizieren die nachgelagerte Kette.}

\chapter{Die zwei Kontinuumsketten}
\label{ch:two-chains}

Die beiden Ketten teilen sich ein Substrat (die Q-Ereignisse), verwenden aber
unterschiedliche Operatoren darauf. Dieses Kapitel führt beide Grenzwerte vor und
formuliert dann die Signatur-Inkompatibilität, die sie trennt.

\section{$L_Q \to -\nabla^2$: Der flache Laplace-Operator}
\label{ch:flat-laplacian}

\begin{theorem}[Diskretes Spektrum des Ketten-Laplace-Operators]
Der Graphen-Laplace-Operator einer 1D-Kette aus $N$ Plätzen hat das exakte Spektrum
\begin{equation}
\lambda_k = \frac{1}{\Delta x^2}\left[2 - 2\cos\frac{\pi k}{N+1}\right],
\qquad k=1,\dots,N,
\end{equation}
mit $\Delta x$ als Gitterabstand. Im Kontinuumsgrenzwert $N\to\infty$, $\Delta x\to0$
bei fester Länge gilt
\begin{equation}
\lambda_k \to (\pi k)^2 .
\end{equation}
\end{theorem}

\begin{proof}
Der Ketten-Laplace-Operator ist die tridiagonale Toeplitz-Matrix mit Diagonale
$2/\Delta x^2$ und Nebendiagonale $-1/\Delta x^2$; ihre Eigenvektoren sind die diskreten
Sinusmoden $\sin(\pi k j/(N+1))$.
\end{proof}

\textbf{Verifikation.} \texttt{GraphLaplacianContinuumTests} baut Ketten mit
$N=32,64,128,256$, diagonalisiert $L_Q$ mittels MathNet.Numerics und vergleicht mit der
geschlossenen Form. Der relative Fehler \emph{sinkt} mit $N$ (Konvergenz zweiter Ordnung,
Fehler $\sim4\times$ pro Verdopplung), womit festgestellt ist, dass $L_Q$ mit
kontrolliertem Fehler gegen den flachen Laplace-Operator konvergiert.
\textbf{Ergebnis: BESTANDEN.}

\theoremnotes
{Stehende Wellen auf einer Kette sind diskrete Sinusmoden; nimmt man $N\to\infty$ bei
fester Länge, wird ihre Dispersionsrelation das übliche Impulsquadrat-Spektrum eines
freien Teilchens auf einer Linie.}
{Dies ist die quantitative Brücke von der diskreten Ontologie zum Kontinuum: eine exakte
geschlossene Form, keine numerische Anpassung, und ihre Konvergenz ist bis zur zweiten
Ordnung kontrolliert.}
{Sie ist eindimensional, ungewichtet und nur elliptisch. Sie sagt nichts über die
Lorentz-Kette oder über Krümmung.}

\historynotes
{Die Frage war, ob die diskrete Ontologie überhaupt einen kontrollierten
Kontinuumsgrenzwert besitzt --- der Einwand des Gutachters ``kein kontrollierter
Kontinuumsgrenzwert''.}
{Ein kontrollierter Grenzwert trennt ein diskretes Spielzeugmodell von einem Kandidaten
für Kontinuumsphysik; ohne ihn hat das Schr\"odinger-Programm keine Brücke zum Raum.}
{Naive Finite-Differenzen-Grenzwerte waren bekannt, aber ein Spektrum in geschlossener
Form mit nachgewiesener Konvergenz zweiter Ordnung ($\sim4\times$/Verdopplung) war für
$L_Q$ noch nicht etabliert.}
{Akzeptiert, weil es eine exakte geschlossene Form (keine Anpassung) ist und der Fehler in
\texttt{GraphLaplacianContinuumTests} nachweislich mit $N$ sinkt.}

\section{BDG $\to \square$: Der d'Alembert-Operator}
\label{ch:dalembertian}

Der Kausalmengen-d'Alembert-Operator ist der Benincasa--Dowker--Glaser-(BDG-)Operator:
eine gerichtete, alternierende Summe über kausale Intervalle der Streuung.

\begin{theorem}[BDG-Konvergenz]
Der BDG-Operator konvergiert mit wachsender Streudichte gegen den Kontinuums-d'Alembert-
Operator $\square$, mit führendem Abbruchfehler $O(h^2)$.
\end{theorem}

\textbf{Verifikation.} \texttt{BDGOperatorContinuumTests} verifiziert die Konvergenz mit
einer Toleranz $10^{-2}$ (der führende $O(h^2)$-Fehler bei $h=1/16$ ist
$\approx4.37\times10^{-3}$); der Fehler fällt $\sim4\times$ pro $h$-Halbierung.
\textbf{Ergebnis: BESTANDEN.}

\theoremnotes
{Ein diskretisierter Wellenoperator auf einer Kausalmenge muss alternierend über
Vergangenheits- und Zukunftsintervalle summieren, um die führenden Terme zu kürzen; die
BDG-Vorschrift tut genau dies und lässt im dichten Grenzfall den
Kontinuums-d'Alembert-Operator zurück.}
{Sie etabliert die Lorentz-Hälfte des Kontinuumsprogramms mit kontrolliertem
$O(h^2)$-Fehler, passend zum Kontrollniveau der elliptischen Hälfte.}
{Der Operator ist aus der Kausalmengentheorie importiert, nicht aus den vier Primitiven
abgeleitet, und die Toleranz ist führender Ordnung ($10^{-2}$), sodass nur die niedrigsten
nichttrivialen Terme geprüft werden.}

\section{Die Quanten-Gravitations-Brücke}
\label{ch:bridge}

Die beiden Kontinuumsketten sind \emph{partiell verbunden} — und ihre Signaturen sind
inkompatibel.

\begin{proposition}[Zwei verschiedene Operatoren]
\begin{enumerate}
\item $L_Q$ ist ein \emph{ungerichteter, positiver} Graphen-Laplace-Operator, dessen
Kontinuumsgrenzwert $-\nabla^2$ ist (Riemannsch, elliptisch, positiv semidefinit).
\item Der BDG-Operator ist ein \emph{gerichteter, alternierender} Kausalmengen-Operator,
dessen Kontinuumsgrenzwert $\square$ ist (Lorentz-artig, hyperbolisch, indefinit).
\end{enumerate}
\end{proposition}

\textbf{Verifikation.} \texttt{QuantumGravityBridgeTests} (drei Tests) stellen fest:
\begin{enumerate}
\item $L_Q$ ist positiv semidefinit (minimaler Eigenwert $\ge0$);
\item der BDG-Operator ist indefinit (besitzt negative Eigenwerte);
\item die beiden sind in der Signatur inkompatibel — der Graphen-Laplace-Operator kann
nicht als Lorentz-d'Alembert-Operator dienen (BdgUniquenessAnalyzer O3/O6 verwirft ihn
explizit).
\end{enumerate}
\textbf{Ergebnis: BESTANDEN (3/3).}

\begin{remark}
Dies ist die wichtigste strukturelle Tatsache des Kontinuumsprogramms: AT besitzt
\emph{zwei} Kettentypen, einen Riemannschen ($L_Q\to-\nabla^2$) und einen
Lorentz-artigen ($\mathrm{BDG}\to\square$), die in der Signatur \emph{disjunkt} sind. Die
Brücke zwischen ihnen ist nativ noch nicht geschlossen.
\end{remark}

\theoremnotes
{Ein positiv semidefiniter Operator kann nicht indefinit sein, sodass der elliptische und
der hyperbolische Kontinuumsgrenzwert nicht dasselbe Objekt sein können; die Q-Ereignisse
unterbestimmen, welcher Operator zu verwenden ist.}
{Diese Tatsache macht die Einstein-Wiederherstellung \emph{logisch} statt mathematisch:
Die Schr\"odinger-Kette und die Einstein-Kette sind getrennt, und ihre Verbindung ist das
offene Problem.}
{Sie formuliert Inkompatibilität, keine Konstruktion; sie ist ein Kein-Brücke-Ergebnis,
keine Brücke.}

%=============================================================================

\chapterend
{$L_Q\to-\nabla^2$ (exakt, zweiter Ordnung) und $\mathrm{BDG}\to\square$ ($O(h^2)$) sind
beide kontrolliert; die beiden Ketten sind in der Signatur disjunkt.}
{Keine native Brücke verbindet die elliptische und die Lorentz-Kette.}
{\texttt{GraphLaplacianContinuumTests}, \texttt{BDGOperatorContinuumTests},
\texttt{QuantumGravityBridgeTests} (3) --- alle BESTANDEN.}

\chapter{Der gewichtete Laplace-Operator, Laplace--Beltrami und gekrümmte Schr\"odinger}
\label{ch:weighted}

Nachdem die flache und die Lorentz-Kette etabliert sind, wendet sich dieses Kapitel der
Krümmung zu. Der gewichtete Laplace-Operator $L_W$ ist die natürliche gekrümmte
Verallgemeinerung von $L_Q$: Er liefert eine wohldefinierte gekrümmte
Schr\"odinger-Gleichung und eine Laplace--Beltrami-Näherung, obwohl die Metrik, die ihn
zu einem echten $\Delta_g$ machen würde, extern bleibt.

\section{Der gewichtete Laplace-Operator}

\begin{definition}[Gewichteter Laplace-Operator]
Gegeben die räumliche Kopplungsmatrix $K_{ij}$, ist der gewichtete
Graphen-Laplace-Operator
\begin{equation}
L_W = D_K - K,
\qquad (D_K)_{ii} = \sum_j K_{ij}.
\end{equation}
Dies ist der übliche \emph{unnormalisierte} Graphen-Laplace-Operator (Belkin--Niyogi),
gültig bei uniformer Dichte.
\end{definition}

\begin{proposition}[Eigenschaften]
$L_W$ ist symmetrisch, hat Null-Zeilensumme, ist positiv semidefinit und reduziert sich
auf den ungewichteten Laplace-Operator, wenn $K=A$.
\end{proposition}

\textbf{Verifikation.} \texttt{WeightedLaplacianTests} (vier Tests) bestätigen Symmetrie,
Null-Zeilensumme, positive Semidefinitheit und Reduktion auf den ungewichteten Fall.
\textbf{Ergebnis: BESTANDEN (4/4).}

\theoremnotes
{Die Gewichtung der Kanten des Graphen erlaubt es, die Geometrie in den Kopplungen zu
kodieren; der unnormalisierte Laplace-Operator ist der Operator, dessen
Nullenergie-Zustand auf dem gewichteten Graphen uniform ist.}
{$L_W$ ist das Werkzeug, das die Theorie von der flachen Kette in den gekrümmten Raum
trägt, und seine Grundeigenschaften sind verifiziert statt angenommen.}
{Der unnormalisierte Laplace-Operator ist nur bei uniformer Dichte der richtige Operator;
bei nicht-uniformer Dichte weichen die normalisierten Varianten ab, und nichts davon
liefert die Metrik selbst.}

\section{Die Laplace--Beltrami-Näherung}
\label{ch:laplace-beltrami}

\begin{theorem}[Laplace--Beltrami aus $L_W$]
$L_W$ approximiert den Laplace--Beltrami-Operator $\Delta_g$ im flachen Grenzfall und
reproduziert übliche Mannigfaltigkeitsbeispiele.
\end{theorem}

\textbf{Verifikation.} \texttt{LaplaceBeltramiTests} (drei Tests):
\begin{enumerate}
\item \emph{Flacher Grenzfall}: $L_W$ erhält das flache Spektrum;
\item \emph{Variable Gewichte}: nicht-uniformes $K$ ändert das Spektrum;
\item \emph{Bekannte Mannigfaltigkeit}: der $S^1$-Kreisgraph reproduziert das
Kreisspektrum.
\end{enumerate}
Eine subtile Entartung wurde gefunden und per Audit behoben: Die Nichtnull-$S^1$-Moden
sind zweifach entartet ($e^{\pm ik\theta}$), sodass in aufsteigender Reihenfolge das
erste Auftreten der Mode $k$ beim Index $2k-1$ liegt, nicht bei $k$.
\textbf{Ergebnis: BESTANDEN (3/3).}

\theoremnotes
{Der Laplace--Beltrami-Operator ist ``der Laplace-Operator mit der Metrik darin''; ihn aus
einem gewichteten Graphen zu approximieren bedeutet, das Spektrum der Mannigfaltigkeit
aus den Kopplungen wiederzugewinnen.}
{Er verbindet den Graphenoperator mit dem üblichen geometrischen Operator des gekrümmten
Raums und legt die $S^1$-Entartungsfeinheit offen (und behebt sie).}
{Er ist eine Näherung, kein Satz für beliebiges $g$, und die Metrik, die $L_W$ zu einem
echten $\Delta_g$ macht, ist weiterhin importiert statt erzeugt.}

\section{Die gekrümmte Schr\"odinger-Gleichung}
\label{ch:curved-schrodinger}

\begin{theorem}[Gekrümmte Schr\"odinger-Gleichung]
Der Operator $L_W$ definiert eine gekrümmte Schr\"odinger-Gleichung
\begin{equation}
i\,\partial_t \psi = L_W \psi .
\end{equation}
Weil $L_W$ selbstadjungiert ist, ist die Evolution unitär und die Norm erhalten:
\begin{equation}
\partial_t \|\psi\|^2 = 0 .
\end{equation}
\end{theorem}

\textbf{Verifikation.} \texttt{CurvedSchrodingerTests} (drei Tests): Unitarität
(Normerhaltung), Reduktion auf die flache Schr\"odinger-Gleichung und die Existenz eines
gekrümmten Operators. \textbf{Ergebnis: BESTANDEN (3/3).}

\begin{remark}
Dies etabliert die \emph{Schr\"odinger}-Seite der Brücke in den gekrümmten Raum: Der
gewichtete Laplace-Operator $L_W$ erzeugt eine wohldefinierte, unitäre gekrümmte
Schr\"odinger-Gleichung. Die \emph{Metrik}, die $L_W$ zu einem echten $\Delta_g$ machen
würde, bleibt extern (Kapitel~\ref{ch:metric-emergence}).
\end{remark}

\theoremnotes
{Selbstadjungiertheit ist genau das, was Unitarität garantiert, sodass die gekrümmte
Gleichung die Normerhaltung gratis erbt.}
{Sie schließt die Schr\"odinger-Seite des Programms für den gekrümmten Raum: Eine
wohldefinierte, getestete, unitäre gekrümmte Dynamik existiert.}
{Die gekrümmte Gleichung ist formal, bis eine Metrik geliefert wird; die Metrikseite ist
extern, was der Grund dafür ist, dass die Einstein-Seite partiell bleibt.}

%=============================================================================

\chapterend
{$L_W$ ist symmetrisch/PSD/nullsummig; er liefert eine Laplace--Beltrami-Näherung und eine
unitäre gekrümmte Schr\"odinger-Gleichung.}
{Die Metrik, die $L_W$ zu einem echten $\Delta_g$ machen würde, bleibt extern.}
{\texttt{WeightedLaplacianTests} (4), \texttt{LaplaceBeltramiTests} (3),
\texttt{CurvedSchrodingerTests} (3) --- alle BESTANDEN.}

\chapter{Die Einstein-Tensor-Kette}
\label{ch:einstein-tensor}

Dieses Kapitel führt die übliche differentialgeometrische Kette von einer Metrik zum
Einstein-Tensor vor. Die Kette ist voll funktionsfähig und getestet; der Sinn ihrer
Aufnahme besteht darin zu zeigen, \emph{genau wo} sie bricht, wenn man versucht, ihr ATs
eigene Objekte zu füttern.

\section{Standard-Differentialgeometrie}

\begin{definition}[Christoffel-Symbole]
\begin{equation}
\Gamma^\lambda_{\mu\nu} = \tfrac12\, g^{\lambda\sigma}
\left( \partial_\mu g_{\sigma\nu} + \partial_\nu g_{\sigma\mu} - \partial_\sigma g_{\mu\nu} \right).
\end{equation}
\end{definition}

\begin{definition}[Riemann-Krümmung]
\begin{equation}
R^\rho_{\ \sigma\mu\nu} = \partial_\mu \Gamma^\rho_{\nu\sigma}
- \partial_\nu \Gamma^\rho_{\mu\sigma}
+ \Gamma^\rho_{\mu\lambda}\Gamma^\lambda_{\nu\sigma}
- \Gamma^\rho_{\nu\lambda}\Gamma^\lambda_{\mu\sigma}.
\end{equation}
\end{definition}

\begin{definition}[Ricci- und Einstein-Tensoren]
\begin{equation}
R_{\sigma\nu} = R^\rho_{\ \sigma\rho\nu}, \qquad
G_{\mu\nu} = R_{\mu\nu} - \tfrac12 R\, g_{\mu\nu}.
\end{equation}
\end{definition}

\section{Verifizierte Beispiele}

\textbf{Einheits-2-Sphäre} $g=\mathrm{diag}(1,\sin^2\theta)$ bei $\theta=\pi/4$:
\begin{equation}
\Gamma^\theta_{\phi\phi} = -\sin\theta\cos\theta = -0.5,
\qquad K = 1, \qquad R = 2K = 2, \qquad R_{\mu\nu}=K g_{\mu\nu}.
\end{equation}

\textbf{Einheits-3-Sphäre} $g=\mathrm{diag}(1,\sin^2\chi,\sin^2\chi\sin^2\theta)$ bei
$\chi=\theta=\pi/4$: $R_{\mu\nu}=2g_{\mu\nu}$, $R=6$, daher
\begin{equation}
G_{\mu\nu} = R_{\mu\nu} - \tfrac12 R g_{\mu\nu} = -g_{\mu\nu}
= -\mathrm{diag}(1,\tfrac12,\tfrac14).
\end{equation}

\begin{theorem}[2D-Einstein-Tensor verschwindet]
In zwei Dimensionen gilt $G_{\mu\nu}\equiv0$ identisch
($R_{\mu\nu}=\tfrac12 R g_{\mu\nu}$); ein nichttrivialer $G_{\mu\nu}$ erfordert Dimension
$\ge3$.
\end{theorem}

\textbf{Verifikation.} \texttt{EinsteinTensorTests} (vier Tests:
Metrik$\to$Christoffel, Christoffel$\to$Riemann, Riemann$\to$Ricci,
Ricci$\to$Einstein) und \texttt{EinsteinTensorIntegrationTests} (vier Tests auf dem
integrierten Builder \texttt{EinsteinTensorBuilder}). \textbf{Ergebnis: BESTANDEN (8/8).}

\theoremnotes
{In zwei Dimensionen ist der Ricci-Tensor gezwungen, proportional zur Metrik zu sein,
sodass der spurfreie Einstein-Tensor identisch verschwindet; Krümmungsdynamik braucht
mindestens drei Dimensionen, um Platz zu haben.}
{Es erklärt innerhalb der Standardgeometrie, warum eine dynamische Einstein-Gleichung in
$d\ge3$ lebt --- eine Tatsache, die das Dimensionalitätsprogramm
(Kapitel~\ref{ch:complexity}) in ein Selektionsargument verwandelt.}
{Es ist Standard-Differentialgeometrie, keine AT-Ableitung; es erzeugt für sich genommen
kein $3{+}1$.}

\begin{remark}[Wo die Kette in AT bricht]
Die Standardkette ist voll funktionsfähig. Aber ATs eigene Analysatoren
(\texttt{QuantumGravityEmergenceAnalyzer}, \texttt{GrBridgeAnalyzer}) \emph{beschreiben}
die Kette als Zeichenketten (\texttt{GeoStep}) und \emph{berechnen} sie nie. Die Kette
bricht bei Schritt~1 — \emph{Metrik} $\to$ \emph{Christoffels} —, weil AT kein natives
Metrikfeld $g_{\mu\nu}(x)$ besitzt: Es ist importiert
(Kapitel~\ref{ch:metric-emergence}).
\end{remark}

%=============================================================================

\chapter{Metrikentstehung und Ursprungsabschluss}
\label{ch:metric-emergence}

Dieses Kapitel fragt, wie eine Metrik überhaupt entstehen kann, und hält die viergliedrige
Kette fest, die den Metrik\emph{ursprung} schließt --- während die Metrik\emph{dynamik}
importiert bleibt.

\section{Die kausale Distanzstruktur}

\begin{definition}[Kausales Volumen und Distanz]
Für ein kausales Intervall der Tiefe $D$ in $d$ Dimensionen ist das kausale Volumen
\begin{equation}
N(D) \propto D^d,
\end{equation}
und die kausale Distanz ist
\begin{equation}
L = N^{1/d}.
\end{equation}
\end{definition}

\textbf{Verifikation.} \texttt{MetricGenerationTests} und \texttt{MetricEmergenceTests}
bestätigen: $L=N^{1/d}$ reproduziert die Tiefe exakt ($d=2,3,4$; $D=1,\dots,8$); die
Dimension wird aus dem Volumenwachstum wiedergewonnen; und die resultierende
Distanzmatrix erfüllt alle vier Metrik-Axiome (Nichtnegativität, Identität der
Ununterscheidbaren, Symmetrie, Dreiecksungleichung). \textbf{Ergebnis: BESTANDEN.}

\section{Der konforme Faktor aus dem Zählmaß}

\begin{theorem}[Konformer Faktor]
Das Zählmaß legt das Volumenelement $\sqrt{|g|}=\rho$ (Dichte $\rho$) fest. Für einen
konform flachen Ansatz $g=f\,\eta$ gilt
\begin{equation}
\sqrt{|g|} = f^{\,d/2} = \rho \quad\Longrightarrow\quad f = \rho^{\,2/d}.
\end{equation}

\eqnotes
{Die Ereignisdichte ist das natürliche Volumenelement, und in einem konform flachen
Ansatz legt das Volumenelement den konformen Faktor fest.}
{Das Zählmaß bestimmt den einen verbleibenden skalaren Freiheitsgrad.}
{Es schließt den konformen \emph{Faktor} nativ, die native Hälfte der
Metrik-Ursprungskette.}
{Nur konform flach; eine allgemeine Metrik wird nativ nicht erzeugt.}
\end{theorem}

\textbf{Verifikation.} \texttt{MetricEmergenceTests} bestätigt den Rundlauf
$f^{d/2}=\rho$ für $\rho\in\{0.5,1,2,8\}$, $d\in\{2,3,4\}$. Der konform flache Kandidat
$g=f\eta$ ist konstruierbar: konstantes $f$ liefert $R=0$; $f=1+\tfrac12 x^2$ liefert
$R=-0.6145$, passend zum Liouville-Ergebnis $K=-\tfrac1{2f}\nabla^2\ln f$.
\textbf{Ergebnis: BESTANDEN.}

\theoremnotes
{Die Dichte der Ereignisse ist der natürliche Kandidat für das Volumenelement, und in
einem konform flachen Ansatz legt allein das Volumenelement den konformen Faktor fest.}
{Dies schließt den konformen \emph{Faktor} nativ: Das Zählmaß bestimmt den einzigen
verbleibenden skalaren Freiheitsgrad.}
{Es setzt einen konform flachen Ansatz voraus; eine allgemeine Metrik wird nativ nicht
erzeugt, nur ihre konforme Klasse plus ein flacher Repräsentant.}

\section{Die konforme Klasse: Satz von Malament}

\begin{theorem}[Malament 1977; Hawking--King--McCarthy 1976]
Eine Bijektion, die die Kausalordnung $\prec$ zwischen zwei
vergangenheits-und-zukunfts-unterscheidenden Raumzeiten erhält, ist eine konforme
Isometrie: Die Kausalordnung bestimmt die Metrik \emph{bis auf einen konformen Faktor}.
\end{theorem}

\textbf{Verifikation.} \texttt{ConformalStructureTests} und \texttt{MetricOriginTests}
stellen fest: Die Kausalordnung ist eine gültige Halbordnung, deren Rand der Lichtkegel
ist; die Nullstruktur ist unter $g\to f g$ ($f>0$) invariant, wird aber durch eine
nicht-konforme Umskalierung $g=\mathrm{diag}(-1,2)$ \emph{verändert}; daher greift der
Lichtkegel \emph{exakt} die konforme Klasse heraus. \textbf{Ergebnis: BESTANDEN.}

\theoremnotes
{Zwei Beobachter, die sich über ``was vor was kommt'' einig sind, sind sich über den
Lichtkegel einig, mithin über die Metrik bis auf eine lokale Gesamt-Umskalierung --- die
konforme Klasse.}
{Dies ist der importierte Satz, der die Kausalordnung in Geometrie verwandelt; einen
bewiesenen Satz zu importieren ist übliche Praxis, kein Mangel.}
{Er ist importiert, nicht AT-abgeleitet, und benötigt
vergangenheits-und-zukunfts-unterscheidende Raumzeiten; er legt nur die konforme Klasse
fest und lässt den Faktor separat zu liefern.}

\begin{theorem}[Abschluss des Metrikursprungs]
Die konforme-Klassen-``Lücke'' ist ein \emph{importierter Satz} (Malament), keine
Theorielücke. Der Metrikursprung ist geschlossen:
\begin{equation}
Q\text{-Ereignisse} \to \text{Kausalordnung (nativ)}
\to \text{konforme Klasse (bewiesen)}
\to \text{konformer Faktor (nativ)}
\to g_{\mu\nu}.
\end{equation}
\end{theorem}

\textbf{Klassifikation des Metrikursprungs:}

\begin{center}
\begin{tabular}{lll}
\toprule
Glied & Status & Natur \\
\midrule
Q-Ereignisse & \textbf{NATIV} & AT-abgeleitetes Primitiv \\
Kausalordnung & \textbf{NATIV} & Präzedenz aus der Individuation \\
Konforme Klasse & \textbf{IMPORTIERT} & Malament/HKM — \emph{bewiesen} \\
Konformer Faktor & \textbf{NATIV} & Zählmaß, $f=\rho^{2/d}$ \\
Volles $g_{\mu\nu}$ & \textbf{bestimmt} & Klasse $\times$ Faktor \\
\bottomrule
\end{tabular}
\end{center}

\begin{remark}
Der Metrikursprung ist \emph{bereits gelöst}: Einen bewiesenen Satz zu importieren ist
üblich, kein Mangel. Der wirklich offene Punkt ist die \emph{native} Metrik $\to$
Operator-Kopplung (die ``G4''-Lücke): AT erzeugt $g_{\mu\nu}$ nicht aus Q-Ereignissen,
es importiert es.
\end{remark}

\theoremnotes
{Die Metrik ist aus drei nativen Gliedern und einem bewiesenen Import zusammengesetzt;
die ``Lücke'' ist ein bewiesener Satz, sodass der Ursprung wirklich geschlossen ist.}
{Sie ist die Antwort der Theorie auf ``woher kommt die Geometrie'': eine Kette, kein
Postulat.}
{Den Ursprung zu schließen schließt nicht die Dynamik: Die native
Metrik-Operator-Kopplung (G4) bleibt offen, weshalb die Einstein-Wiederherstellung
logisch statt mathematisch ist.}

\historynotes
{Die Frage war, ob eine Metrik $g_{\mu\nu}$ überhaupt aus Q-Ereignissen entstehen kann und
wo die Theorie anhalten muss.}
{Geometrie ist das Substrat der Einstein-Seite; ein nativer Ursprung würde die größte
verbleibende Lücke im Kontinuumsprogramm schließen.}
{Direkte Konstruktion aus Q-Ereignissen scheiterte (der metrikabhängige Operator fehlt);
der Weg, der funktionierte, importiert den Satz von Malament für die konforme Klasse und
liefert den konformen Faktor nativ.}
{Akzeptiert als Abschluss des Metrik\emph{ursprungs} --- der Import ist ein bewiesener
Satz, keine Lücke ---, während die Metrik\emph{dynamik} (G4) ehrlich importiert bleibt.}

%=============================================================================
\chapterend
{Der Metrikursprung ist geschlossen: Q-Ereignisse $\to$ Kausalordnung $\to$ konforme
Klasse (Malament, bewiesen) $\to$ konformer Faktor ($f=\rho^{2/d}$, nativ) $\to$ Metrik.}
{Die native Metrik-Operator-Kopplung (G4) bleibt offen; die Metrikdynamik ist importiert.}
{\texttt{MetricGenerationTests} (4), \texttt{MetricEmergenceTests} (4),
\texttt{ConformalStructureTests} (3), \texttt{MetricOriginTests} (3) --- alle BESTANDEN.}

\part{Die Ableitungshierarchie}
%=============================================================================

Dieser Teil führt die Ableitungshierarchie von den Kontinuumsobjekten hinab zur konkreten
Physik: Dimensionalität, Eichstruktur, Flavour, Gravitation und Kosmologie. Jedes Kapitel
hält fest, was \derived, was \realunder und was \drawn ist, mit den
Konfidenzeinschätzungen, die die Geschichte der feindlichen Begutachtung verfeinert hat.

\chapter{Komplexität und räumliche Dimensionalität}
\label{ch:complexity}

\section{Das Komplexitätsfunktional}

Das Komplexitätsargument ist eine \textbf{gewichtete Sechs-Komponenten-Zerlegung} der
Bedingungen für Beobachter (\texttt{ComplexityOptimumAnalyzer.cs}), kein einzelnes
skalares Variationsfunktional.

\begin{center}
\begin{tabular}{lcc}
\toprule
Komponente & Gewicht & $M^2$-Fenster \\
\midrule
Strukturbildung & 3,0 & 2--8 \\
Teilchenstabilität & 2,5 & 3--7 \\
Chemiepotenzial & 4,0 & 3--5 \\
Informationskapazität & 2,0 & 3--10 \\
Evolutionspotenzial & 1,5 & 2--6 \\
Beobachterviabilität & 5,0 & 4--6 \\
\bottomrule
\end{tabular}
\end{center}

\section{Räumliche Dimensionalität $d=3+1$}

\begin{theorem}[Eindeutige Dimensionalität]
$d=3+1$ ist der \emph{eindeutige} Wert, der gleichzeitig erfüllt: den Satz von Bertrand
(geschlossene Bahnen nur für $1/r^2$), das Gaußsche Gesetz ($F\propto1/r^{d-1}$, Coulomb
nur in 3D), Knotenstabilität (Kodimension 2), Huygens (scharfe Ausbreitung in ungeradem
$d$) und 2 GR-Polarisationen. $d=2+1$ lässt Gravitation/Atome/Knoten scheitern;
$d\ge4+1$ lässt Bahnen und Knoten scheitern.
\end{theorem}

\begin{remark}
$M^2\approx5$ ist der Beobachterviabilitäts-\emph{Schnitt} der Komponentenfenster, keine
von Hand eingesetzte Zahl. Konfidenz T-02 = 0,85; dies ist ein
Fensterschnitt-Argument, kein Variationssatz. Der Analysator hält fest, dass die
Generationszahl $G\approx3$ ein \emph{Plateau} ist, kein Maximum.
\end{remark}

\theoremnotes
{Jede Anforderung --- geschlossene Bahnen, Coulomb-Kräfte, Knoten, scharfe Signale,
Gravitationspolarisationen --- ist ein Filter auf die Dimension, und ihr Schnitt ist ein
einziger Wert.}
{Die Dimensionalität ist abgeleitet statt angenommen, womit eine der größten
``Warum''-Fragen mit einem testartigen Schnittargument geschlossen wird.}
{Es ist ein Fensterschnitt-Argument (Konfidenz 0,85), kein Variationssatz, und es leitet
nur die \emph{Raumzeit}-3 ab; die \emph{innere} 3 (Generationen, Farbe) ist eine
separate, selektierte Größe (Kapitel~\ref{ch:nogo}).}

%=============================================================================

\chapterend
{$d=3{+}1$ ist der eindeutige Wert, der Bertrand, Gauß, Knotenstabilität, Huygens und die
zwei GR-Polarisationen gleichzeitig erfüllt; $M^2\approx5$ ist das
Beobachterviabilitäts-Fenster.}
{Das Argument ist ein Fensterschnitt, kein Variationssatz (T-02 $=$ 0,85).}
{Analyse via \texttt{ComplexityOptimumAnalyzer}; nicht Teil des ausführbaren
Testinventars.}

\chapter{Eichstruktur}
\label{ch:gauge}

Dieses Kapitel hält das Ergebnis des Eichprogramms fest: $U(1)$ abgeleitet, $SU(2)$
emergent, $SU(3)$ kontingent. Die Struktur ist ableitbar; die Anzahl ist es nicht.

\section{$U(1)$ — abgeleitet}

\begin{theorem}[$U(1)$ aus dem Kreis]
Die Phase lebt auf $S^1$. Ihre Isometriegruppe ist
\begin{equation}
\mathrm{Aut}(S^1) = U(1), \qquad \pi_1(S^1) = \mathbb{Z},
\end{equation}

\eqnotes
{Die Phase ist ein Winkel, lebt also auf einem Kreis; die Symmetrien des Kreises sind
$U(1)$, und ihn ganzzahlig oft zu umrunden ist eine erhaltene Ladung.}
{$U(1)$ ist die Symmetrie der Phase selbst.}
{Sie entfernt eine freie Eichannahme aus dem Standardwerkzeugkasten.}
{Nur Struktur: Die Maxwell-Wirkung ist weiterhin entlehnt, nicht erzeugt.}
und die Windungszahl $\oint\nabla\theta\cdot dl = 2\pi n$ liefert ganzzahlige topologische
Ladung. (Konfidenz 0,95.)
\end{theorem}

\theoremnotes
{Die Phase ist ein Winkel, lebt also auf einem Kreis; die Symmetrien eines Kreises sind
$U(1)$, und den Kreis ganzzahlig oft zu umrunden ist eine erhaltene Ladung.}
{Dies ist das sauberste Ergebnis des Eichprogramms: $U(1)$ ist nicht gewählt, sondern die
Symmetrie der Phase selbst, zu einem Satz erhoben (Konfidenz 0,95).}
{Es leitet die Gruppenstruktur ab, nicht die Eichdynamik: Die Maxwell-Wirkung ist
entlehnt, nicht erzeugt.}

\historynotes
{Die Frage war, ob die $U(1)$-Elektromagnetik eine separate Eingabe ist oder aus der
Ontologie der Phase folgt.}
{$U(1)$ ist die älteste Eichsymmetrie; sie aus dem Kreis abzuleiten würde eine freie
Annahme aus dem Standardwerkzeugkasten entfernen.}
{Gruppentheoretische Zugänge \emph{postulieren} $U(1)$; der Defekt-/Windungsweg
versuchte, sie aus $\mathrm{Aut}(S^1)$ und $\pi_1(S^1)=\mathbb Z$ \emph{abzuleiten}.}
{Akzeptiert als das sauberste Ergebnis des Eichprogramms (Konfidenz 0,95): keine Wahl,
sondern die Symmetrie der Phase selbst --- wobei zugestanden wird, dass die
Maxwell-Wirkung weiterhin entlehnt ist.}

\section{$SU(2)$ und $SU(3)$ — reale, nicht abgeleitete Struktur}

\begin{center}
\begin{tabular}{lll}
\toprule
Faktor & Status & Konfidenz \\
\midrule
$U(1)$ & \derived & 0,95 \\
$SU(2)$ & \realunder\ (emergent) & 0,70 \\
$SU(3)$-Struktur & \realunder\ (emergent) & 0,10 \\
Farbanzahl $n=3$ & \drawn & — \\
\bottomrule
\end{tabular}
\end{center}

Binäre Windung $\{n=\pm1\}$ liefert das minimale Dublett; die Bloch-Sphäre liefert
$SO(3)=SU(2)/Z_2$; die Spinor-Doppelüberlagerung liefert $SU(2)$. Der Automorphismus des
$n$-Defekt-Moduliraums erfüllt $\mathrm{Aut}(C^n/S_n)\supseteq SU(n)$ und leitet die
$SU(3)$-\emph{Struktur} ab, aber nicht die Anzahl.

\begin{remark}
AT leitet die Eichgruppen-\textbf{Struktur} ab (welche Gruppe), \textbf{nicht} die
Eich\textbf{dynamik} (Maxwell- / Yang--Mills-Wirkungen sind entlehnt; die
8-Gluon-Algebra ist Annahme A-07). Die Anzahl $n=3$ ist nicht festgelegt
(Eichanzahl-No-Go T-09, vorläufig 0,10).
\end{remark}

%=============================================================================

\chapterend
{$U(1)$ ist abgeleitet ($\mathrm{Aut}(S^1)=U(1)$, $\pi_1(S^1)=\mathbb Z$); $SU(2)$ ist
emergent; $SU(3)$ ist kontingent (Struktur, nicht Dynamik).}
{Die Eichanzahl $n=3$ ist nicht abgeleitet (vorläufiges No-Go T-09, 0,10).}
{Mathematisch (Topologie); Konfidenz 0,95 / 0,70 / 0,10; kein eigener ausführbarer Test.}

\chapter{Flavour und die Koide-Relation}
\label{ch:flavor}

Das Flavour-Kapitel ist die klarste einzelne Veranschaulichung der
Struktur/Inhalt-Trennung: Die Operator-\emph{Form} ist abgeleitet, ihr \emph{Spektrum}
ist kontingent, und die Koide-Relation liegt exakt auf der Grenze --- real, präzise,
vorhersagekräftig und nicht ableitbar.

\section{Der Yukawa-Operator}

Der Yukawa-Operator ist der Überlappoperator
\begin{equation}
Y_{ij} = \langle \mathrm{arch}_i \,|\,\mathrm{amplitude}\,|\,\mathrm{arch}_j\rangle .
\end{equation}
Seine \emph{Form} ist abgeleitet; sein \emph{Spektrum} (die Hierarchie $1{:}207{:}3478$)
ist nicht abgeleitet — festgelegt durch die kontingenten Architekturformen.

\section{Die Koide-Relation}

\begin{theorem}[Koide-Relation]
Die Polmassen der geladenen Leptonen erfüllen
\begin{equation}
Q = \frac{\sum m}{(\sum\sqrt{m})^2} = 0.6666605,
\qquad \theta = 44.9997^\circ,
\end{equation}

\eqnotes
{Im $\sqrt m$-(Amplituden-)Raum sitzen die drei geladenen Leptonen unter $45^\circ$ zur
demokratischen Richtung.}
{$Q=2/3$ ist $1/(3\cos^2 45^\circ)$; die Relation handelt von Amplituden, nicht von
Massen.}
{Sie ist der schärfste numerische Hinweis im Flavour-Sektor: real, vorhersagekräftig,
RG-stabil.}
{Ihr Ursprung ist nicht ableitbar (No-Go T-08): real bedeutet nicht abgeleitet.}
äquivalent $Q=2/3$ bei $\theta=\arccos(1/\sqrt2)=45^\circ$. Die Relation ist real
($\sim10^{-5}$ Präzision; Bayes-Faktor vs.\ Zufall $\approx3.2\times10^4$),
vorhersagekräftig (Vorhersage von 1981: $m_\tau=1776.97$ MeV, 1992 bestätigt) und
RG-stabil.
\end{theorem}

\theoremnotes
{Im $\sqrt m$-(d.\,h. Amplituden-)Raum sitzen die drei geladenen Leptonen unter einem
Winkel von $45^\circ$ zur demokratischen Richtung, was exakt die Bedingung $Q=2/3$ ist;
die Relation handelt von Amplituden, nicht von Massen.}
{Sie ist eine reale, vorhersagekräftige, RG-stabile Regelmäßigkeit --- der stärkste
numerische Hinweis im Flavour-Sektor --- und das kanonische Beispiel einer
\realunder-Struktur.}
{Ihr Ursprung ist nicht ableitbar (No-Go T-08 unten): real bedeutet nicht abgeleitet, und
die Relation legt für sich genommen nicht die Drittgenerations-Anzahl fest.}

\historynotes
{Die Koide-Relation steht seit 1981 als unerklärte, hochpräzise Regelmäßigkeit; das
Programm fragte, ob ihr Ursprung aus den Primitiven ableitbar ist.}
{Sie ist der einzige schärfste numerische Hinweis im Flavour-Sektor; ihr Ursprung würde
entweder eine tiefe Struktur bestätigen oder beweisen, dass Struktur bei der Form endet.}
{Sieben Wege wurden versucht --- Symmetrie, Topologie, Attraktoren, Informationsgeometrie,
Gruppentheorie, $m=3$-Abschluss, Moduli --- und alle scheiterten oder waren blockiert.}
{Akzeptiert als \emph{geschlossene Frage} (No-Go T-08): real, vorhersagekräftig,
RG-stabil und nicht ableitbar --- die kanonische \realunder-Struktur, mit dem
falsifizierten Neutrino-Koide-Analogon (T-11).}

\begin{nogo}[Koide-Ursprung — T-08]
Keine Symmetrie ($S_3$), kein Attraktor, keine Topologie ($S^1$) und keine
Informationsgeometrie selektiert $Q=2/3$. (Konfidenz 0,70.)
\end{nogo}

Das Abschluss-Audit (Phase 159) schöpfte sieben Wege aus:

\begin{center}
\begin{tabular}{ll}
\toprule
Weg & Status \\
\midrule
Symmetrie ($S_3$) & No-Go \\
Topologie ($S^1$) & No-Go \\
Attraktoren & No-Go \\
Informationsgeometrie & No-Go \\
Gruppentheorie & No-Go \\
$m=3$-Abschluss & No-Go \\
Moduli-Argumente & Blockiert \\
\bottomrule
\end{tabular}
\end{center}

Der Koide-Ursprung ist eine \emph{geschlossene Frage} — real, vorhersagekräftig,
RG-stabil, \emph{nicht ableitbar} — die kanonische \realunder-Struktur.

\begin{nogo}[Neutrino-Koide — T-11, Falsifikation]
$Q=2/3$ für Neutrinos zu verlangen liefert $Q_{\max}=0.585$ (normal) / $0.500$
(invertiert), beides $<2/3$; die gemessenen $\Delta m^2$ schließen All-Lepton-Koide aus.
(Konfidenz 0,90.)
\end{nogo}

%=============================================================================

\chapterend
{Die Koide-Relation ist real, vorhersagekräftig und RG-stabil; ihr Ursprung ist nicht
ableitbar (T-08); das Neutrino-Koide-Analogon ist falsifiziert (T-11).}
{Keine: Der Koide-Ursprung ist eine geschlossene Frage, und die Falsifikation schloss den
letzten Unterscheider.}
{Berechnungs-Audits ($Q$, $Q_{\max}$); im Koide-Abschluss festgehalten, nicht im
ausführbaren Testinventar.}

\chapter{Gravitation und emergente ART}
\label{ch:gravity}

Gravitation ist in AT ein Phasengradienten-Phänomen. Dieses Kapitel hält den ehrlichen
Status fest: Die Einstein-Wiederherstellung ist \emph{logisch}, keine neue dynamische
Ableitung.

\section{Newtons Konstante}

\begin{proposition}[Dimensionsergebnis]
$G = \ell^2 c^3/\hbar$ ist ein \emph{Einheitenkonsistenz}-Ergebnis, keine Ableitung der
Gravitationsdynamik.
\end{proposition}

\eqnotes
{Das Tripel $(\ell,\tau,\hbar)$ kann genau eine Kombination mit den Dimensionen von $G$
zusammensetzen.}
{$G$ ist durch die Einheiten festgelegt, sobald die Skala $\ell$ festgelegt ist.}
{Es entfernt Newtons Konstante aus der Liste unabhängiger Eingaben.}
{Nur Dimensionsanalyse; es sagt nichts über die Feldgleichungen.}

\theoremnotes
{Das Tripel $(\ell,\tau,\hbar)$ kann genau eine Kombination mit den Dimensionen von $G$
zusammensetzen, sodass Newtons Konstante durch die Einheiten festgelegt ist, sobald die
Skala festgelegt ist.}
{Es entfernt $G$ aus der Liste unabhängiger Konstanten --- die Stärke der Gravitation ist
eine buchhalterische Konsequenz des Skalen-Tripels.}
{Es ist Dimensionsanalyse, keine Dynamik; es sagt nichts über die Feldgleichungen.}

\section{Die Phasengravitations-Kette}

Oszillation $\to$ Phase $\to$ Kausaldichte $\to$ Metrik $\to$ Krümmung $\to$ Einstein.

\begin{theorem}[Wiederherstellung führender Ordnung]
\begin{equation}
G_{\mu\nu} = 8\pi G_{\rm eff}\, T_{\mu\nu} + O(\ell_P^2 R^2),
\end{equation}
mit Planck-Skalen-Korrekturen $\sim10^{-40}$. Newtonsches $1/r^2$, Linseneffekt,
Rotverschiebung, Periheldrehung und zwei Gravitationswellen-Polarisationen mit
Geschwindigkeit $c$ werden reproduziert.
\end{theorem}

Zwei Modifikationen jenseits der ART: (1) zeitveränderliche dunkle Energie
\begin{equation}
\Lambda(t) = \frac{\alpha}{\sqrt{V(t)}},
\end{equation}
wobei $V(t)$ das 4-Volumen des vergangenen Lichtkegels ist; und (2)
Singularitätsauflösung bei $r\sim\ell_P$ (maximale Krümmung $1/\ell_P^2$).

\begin{remark}[Ehrlichkeitshinweis]
Die Kette ist \textbf{logisch, keine neue dynamische Ableitung}: ``gleiche Gleichungen,
gleiche Vorhersagen; kein falsifizierbarer Unterschied'' im schwachen Feld. Der Wert ist
\emph{ontologisch} (die Identifikation des Gravitationspotenzials mit dem Phasenfeld); der
physikalische Gehalt jenseits der ART beschränkt sich auf $\Lambda(t)$ und die
Singularitätsauflösung.
\end{remark}

\theoremnotes
{Wenn das Phasenfeld das Gravitationspotenzial ist, dann wird die übliche Krümmungskette
bis zur führenden Ordnung wiederhergestellt, mit nur Planck-Skalen-Korrekturen ---
Gravitation wird neu etikettiert, nicht ersetzt.}
{Sie zeigt, dass die ontologische Passung im schwachen Feld mit der ART konsistent ist,
während sie die beiden Stellen isoliert, an denen AT wirklich abweicht ($\Lambda(t)$ und
Singularitätsauflösung).}
{Sie ist logisch, nicht mathematisch: Metrik und BDG-Wirkung sind importiert, sodass dies
eine Reinterpretation ist, keine Ableitung der Einstein-Gleichungen.}

\section{Die Radialbeschleunigungsrelation}

\begin{theorem}[Nullparameter-RAR]
\begin{equation}
g_\dagger = \frac{c H_0}{2\pi} \approx 1.05\times10^{-10}\ \mathrm{m/s^2},
\end{equation}

\eqnotes
{$cH_0$ ist die natürliche Beschleunigungsskala einer Kausalmenge; die Division durch
$2\pi$ reproduziert die Milgrom-Skala $a_0$.}
{Sie ist eine Nullparameter-Vorhersage, die auf dem gemessenen Wert landet.}
{Sie ist das schärfste Dunkler-Sektor-Ergebnis der Theorie.}
{Das $2\pi$ ist ein zugestandener numerischer Zufall; dies ist eine Übereinstimmung, noch
keine vollständige Ableitung.}
passend zum gemessenen $a_0$. (Der Faktor $2\pi$ ist ein numerischer Zufall, Phase 144.)
\end{theorem}

\theoremnotes
{Die Kombination $cH_0$ ist die natürliche Beschleunigungsskala einer Kausalmenge; die
Division durch $2\pi$ reproduziert die beobachtete Milgrom-Skala $a_0$.}
{Sie ist eine Nullparameter-Vorhersage, die auf dem gemessenen Wert landet, ein wirklich
scharfes Ergebnis im Dunklen Sektor.}
{Der $2\pi$-Faktor ist ein numerischer Zufall, und die Relation ist eine Übereinstimmung,
noch keine vollständige Ableitung der Dynamik hinter $a_0$.}

%=============================================================================

\chapterend
{$G=\ell^2c^3/\hbar$ ist dimensional; die Einstein-Wiederherstellung führender Ordnung ist
logisch; die RAR $g_\dagger=cH_0/2\pi$ trifft $a_0$ mit null freien Parametern.}
{Native Metrikdynamik (G4); eine eindeutige, scharfe diskriminierende Vorhersage.}
{\texttt{EinsteinTensorTests} (4) und \texttt{EinsteinTensorIntegrationTests} (4)
verifizieren die Standardkette; \texttt{EinsteinRecoveryTests} (3) = TEILWEISE.}

\chapter{Kosmologie}
\label{ch:cosmology}

Die Kosmologie ist das am wenigsten geschlossene Kapitel der Monographie: Sie ist eine
akzeptierte partielle Berechnungsschicht, keine vollendete Ableitung. Der Status wird
ohne Überbeanspruchung festgehalten.

\begin{itemize}
\item \textbf{Expansion} ist eine Interpretation von FLRW (QG-081); jede
Reinterpretation kollabiert zu FLRW (QG-080--089).
\item \textbf{Kausalmengen-$\Lambda$}: $\Lambda\sim1/\sqrt{N}$ ist eine Postdiktion
(Phase 140).
\item \textbf{Pantheon+SH0ES}: AT und $\Lambda$CDM sind beim aktuellen Signal
$w(z)=-1+0.015(1+z)^{3/2}$ nicht unterscheidbar (DATA-001); ein Nachweisbarkeits-Audit
(DATA-002) quantifizierte die für eine Trennung erforderliche Leistung.
\item \textbf{Cluster}: Die dynamische Coma-Masse und der ACCEPT-Röntgengasanteil werden
reproduziert.
\item \textbf{CMB}: eine akzeptierte partielle Berechnungsschicht; vollständiger
$C_\ell$-Löser zurückgestellt (rechnerisch, keine Theorielücke).
\end{itemize}

\begin{remark}
Der Kosmologiestatus ist bewusst partiell: Der Hintergrund und die Kompressionsmaxima des
CMB sind vorhanden, aber das Verdünnungsmaximum und die akustische Phasenverschiebung
erfordern einen Boltzmann-Löser der CAMB-Klasse --- Standard-$\Lambda$CDM-Physik, keine
neue AT-Physik.
\end{remark}

%=============================================================================
\chapterend
{Expansion ist eine FLRW-Reinterpretation; Kausalmengen-$\Lambda\sim1/\sqrt N$ ist eine
Postdiktion; das $w(z)$-Signal ist quantifiziert.}
{Der vollständige CMB-$C_\ell$-Löser ist zurückgestellt (rechnerisch, keine Theorielücke).}
{Partielle Berechnungsschicht (DATA-001/002, CMB-Abschluss-Audits); nicht im ausführbaren
Testinventar.}

\part{Klassifikation, No-Gos und Vorhersagen}
%=============================================================================

\chapter{Die Drei-Kategorien-Taxonomie}
\label{ch:taxonomy}

Die Taxonomie ist das Instrument der ehrlichen Buchführung der Theorie. Sie klassifiziert
\emph{Nichtableitbarkeit}, sodass keine Behauptung als Ableitung durchgehen darf, wenn sie
eine Klassifikation ist.

\begin{definition}[Taxonomie]
\begin{center}
\begin{tabular}{ll}
\toprule
Kategorie & Bedeutung \\
\midrule
\derived & aus den Primitiven per Satz berechenbar \\
\realunder & reale, präzise, vorhersagekräftige Struktur, deren Ursprung nicht berechenbar ist \\
\drawn & eine zufällige Ziehung des Häufigkeitsgesetzes \\
\bottomrule
\end{tabular}
\end{center}
Die Taxonomie klassifiziert \emph{Nichtableitbarkeit}; \realunder\ und \drawn\ werden
\emph{nicht} als Ableitungen beansprucht.
\end{definition}

\section{Die vierzehn klassifizierten Ergebnisse}

\begin{center}
\begin{tabular}{llc}
\toprule
Objekt & Kategorie & Konfidenz \\
\midrule
$U(1)$ & \derived & 0,95 \\
räumliche 3 & \derived & 0,85 \\
$N\ge3$ & \derived & 0,90 \\
Lognormalgesetz & \derived & Satz \\
$SU(2)$ & \realunder\ (emergent) & 0,70 \\
$SU(3)$-Struktur & \realunder\ (emergent) & 0,10 \\
Koide $Q=2/3$ (Realität) & \realunder\ (strukturiert) & 0,90 \\
Koide $45^\circ$ (Ursprung) & \realunder\ (strukturiert) & 0,70 \\
Yukawas / Kopplungen / $\Omega_{DM}$ & \drawn & — \\
$N\le3$ & \drawn & 0,70 \\
Farbanzahl 3 & \drawn & — \\
innere $N=3$ & \derived$\cap$\drawn & 0,70 \\
Neutrino-Koide & FALSIFIZIERT & 0,90 \\
\bottomrule
\end{tabular}
\end{center}

Kategoriekonflikte: \textbf{0}. Zusammengesetzte Objekte: \textbf{2}. Gesamtkonsistenz:
\textbf{0,95}.

%=============================================================================

\chapter{Die No-Go-Sätze}
\label{ch:nogo}

Die No-Go-Sätze sind \emph{bedingt} — relativ zur Keine-neuen-Primitive-Beschränkung und
zum getesteten Wegraum —, keine absoluten logischen Unmöglichkeiten (außer der
Berechnung T-11).

\begin{center}
\begin{tabular}{p{3.4cm}p{9.6cm}c}
\toprule
Satz & Aussage & Konfidenz \\
\midrule
\textbf{T-08} Koide-No-Go & keine Symmetrie/kein Attraktor/keine Topologie/keine Informationsgeometrie selektiert $Q=2/3$ & 0,70 \\
\textbf{T-09} Eichanzahl-No-Go (vorläufig) & kein Prinzip legt $n$ fest; Graphenspektrum/Gittermode ungetestet & 0,10 \\
\textbf{T-10} $N\le3$-No-Go & keine Stabilitäts-/Anomalie-/Darstellungsschranke liefert $N\le3$ & 0,70 \\
\textbf{T-11} Neutrino-Koide-Falsifikation & $Q_{\max}=0.585/0.500 < 2/3$ (Berechnung) & 0,90 \\
\textbf{T-12} Gemeinsame-Kaskade-No-Go & eine vs. drei Kaskaden aus einem Universum untestbar & 0,55 \\
\bottomrule
\end{tabular}
\end{center}

\begin{remark}[Disposition der inneren 3]
Eichanzahl $n=3$ und Multiplizität $N\le3$ sind ein Knoten mit zwei formal unverbundenen
Gesichtern ($N$ vs. $n$; Annahme A-10). Er ist \emph{ungelöst-kontingent} disponiert: Das
Multiplizitätsgesicht ist geschlossen (T-10, 0,70); das Eichanzahl-Gesicht ist
\emph{vorläufig} geschlossen (T-09, 0,10). Dies ist die eine offene Tür der Theorie.
\end{remark}

\theoremnotes
{Jedes No-Go ist eine Aussage darüber, was der getestete Wegraum nicht erreichen kann,
nicht darüber, was logisch unmöglich ist; die ``3'', die die Physik durchzieht, reduziert
sich auf einen Knoten mit zwei Gesichtern.}
{Die No-Gos sind das markanteste Produkt des Programms: Sie lokalisieren die exakte Grenze
zwischen dem Ableitbaren und dem Kontingenten, was der Inhalt der
Struktur/Inhalt-Trennung ist.}
{Sie sind bedingt (relativ zu Keine-neuen-Primitive und den getesteten Wegen), das
Eichanzahl-No-Go T-09 ist nur vorläufig (0,10), und das Gemeinsame-Kaskade-No-Go T-12 ist
schwach (0,55).}

%=============================================================================

\chapterend
{T-08--T-12 sind bedingte No-Gos; der Knoten der inneren 3 ist die eine offene Tür der
Theorie.}
{Das Eichanzahl-Gesicht (T-09, 0,10) ist nur vorläufig geschlossen.}
{T-11 ist eine Berechnung (Falsifikation); T-08 ist Weg-Ausschöpfung; der Rest sind
festgehaltene Audits.}

\chapter{Quantitative Vorhersagen}
\label{ch:predictions}

Der Anspruch der Theorie, wissenschaftlich zu sein, beruht auf der Tatsache, dass sie
Vorhersagen macht, die scheitern können, und bereits eine verfehlt hat. Die Tabelle hält
die lebenden Vorhersagen und die eine falsifizierte Vorhersage fest; die Falsifikation
bleibt in der Tabelle statt begraben zu werden, weil sie der stärkste Beleg dafür ist,
dass das Programm nicht immunisiert ist.

\begin{center}
\begin{tabular}{lll}
\toprule
Vorhersage & Typ & Status \\
\midrule
RAR $g_\dagger=cH_0/(2\pi)\approx1.05\times10^{-10}$ & nullparameterig & trifft $a_0$ \\
$w(z)=-1+0.015(1+z)^{3/2}$ & kosmologisch & Euclid $>3\sigma$ bis $\sim$2030 \\
$\Lambda(t)=\alpha/\sqrt{V(t)}$ & dunkle Energie & Euclid / Roman \\
Lognormales Häufigkeitsgesetz & Verteilungsform & testbar \\
$N\ge3$ (CP-Untergrenze) & a priori & bestätigt \\
Neutrino-Koide $Q=2/3$ & falsifizierbar & \textbf{falsifiziert} \\
\bottomrule
\end{tabular}
\end{center}

Die Theorie ist \emph{falsifizierbar}: Sie machte eine konkrete Vorhersage
(Neutrino-Koide), die falsifiziert wurde, und sie trägt lebende quantitative Vorhersagen,
die sie im \emph{Form}-Sektor von SM $+\Lambda$CDM unterscheiden, obwohl der
\emph{Inhalt} DRAWN ist.

%=============================================================================

\chapter{Umfang und Grenzen}
\label{ch:scope}

Eine Theorie, die ihre Grenzen benennt, ist nützlicher als eine, die sie verbirgt. Die
folgende Liste ist die eigene Buchführung der Theorie darüber, wo sie aufhört; sie ist aus
der Antwort auf die feindliche Begutachtung übernommen und soll als Teil der Theorie
gelesen werden, nicht als Entschuldigung für sie.

\begin{enumerate}
\item Das Eichanzahl-Gesicht des Knotens der inneren 3 ist nur vorläufig geschlossen
(T-09, 0,10).
\item Enzyklopädie-Vollständigkeit $\approx81\%$; das CMB-Kapitel ist eine dokumentierte
partielle Berechnungsschicht.
\item Inhalt ist konstruktionsbedingt kontingent; die Theorie sagt die
Verteilungs\emph{form} voraus, nicht die \emph{Werte}.
\item Die Gemeinsame-Kaskade-Frage ist logisch (nicht empirisch) geschlossen.
\item Die vereinheitlichte Wirkung ist eine Roadmap, kein Ergebnis.
\item Konfidenzzahlen sind unkalibrierte ordinale Ränge (Phase-149/156-Schätzungen).
\item Die Phasengravitations-Kette ist im schwachen Feld ontologisch.
\end{enumerate}

%=============================================================================
\part{Verifikation}
%=============================================================================

\chapter{Das ausführbare Verifikationsprogramm}
\label{ch:verification}

Der Anspruch, dass dies eine \emph{getestete} Theorie ist, beruht auf dem ausführbaren
Programm, das hier zusammengefasst und vollständig in Anhang~\ref{app:tests} inventarisiert
ist. Jedes ``BESTANDEN'' unten ist das Ergebnis eines xUnit-Tests, keine Handberechnung.

\section{Zusammenfassung der verifizierten Ergebnisse}

\begin{center}
\begin{tabular}{lll}
\toprule
Kettenglied & Testdatei & Ergebnis \\
\midrule
$L_Q\to-\nabla^2$ & GraphLaplacianContinuumTests & exaktes Spektrum, 4$\times$/Verdopplung \\
BDG$\to\square$ & BDGOperatorContinuumTests & $O(h^2)$, 4$\times$/Halbierung \\
$L_Q$ vs. BDG-Signatur & QuantumGravityBridgeTests (3) & PSD vs. indefinit \\
kein natives $\Delta_g$ & CurvedSpaceBridgeTests (3) & Brücke fehlt \\
gewichteter Laplace-Operator & WeightedLaplacianTests (4) & symmetrisch/PSD/nullsummig \\
$L_W\to\Delta_g$ & LaplaceBeltramiTests (3) & S¹-Beispiel \\
gekrümmte Schr\"odinger-Gleichung & CurvedSchrodingerTests (3) & unitär \\
Einstein-Wiederherstellung & EinsteinRecoveryTests (3) & TEILWEISE \\
Einstein-Kette & EinsteinTensorTests (4) & Standard 2D/3D \\
Einstein-Integration & EinsteinTensorIntegrationTests (4) & Builder funktioniert \\
Metrik-Erzeugung & MetricGenerationTests (4) & Distanz VORHANDEN \\
Metrikentstehung & MetricEmergenceTests (4) & Faktor nativ \\
konforme Struktur & ConformalStructureTests (3) & Klasse importiert \\
Metrikursprung & MetricOriginTests (3) & Abschluss \\
\bottomrule
\end{tabular}
\end{center}

\section{Die entscheidenden Befunde}

\begin{enumerate}
\item Die \emph{Schr\"odinger}-Seite des Kontinuumsgrenzwerts ist \textbf{kontrolliert und
getestet}: $L_Q\to-\nabla^2\to$ Schr\"odinger (sowie das gekrümmte $L_W$ und die
Laplace--Beltrami-Näherung) sind verifiziert.
\item Die \emph{Einstein}-Seite ist \textbf{teilweise} verifiziert: Die Standardkette
$g\to G_{\mu\nu}$ funktioniert, aber AT hat keine native Metrik (Malament-Import), und
die BDG-Wirkung ist importiert.
\item Die konforme-Klassen-``Lücke'' ist ein \emph{importierter Satz} (Malament), keine
Theorielücke.
\end{enumerate}

%=============================================================================

\chapterend
{Die Schr\"odinger-Seite des Kontinuumsgrenzwerts ist kontrolliert und getestet; die
Einstein-Seite ist partiell; die konforme Klasse ist ein importierter bewiesener Satz.}
{Die native Metrik-Operator-Kopplung (G4).}
{15 Testdateien / 47 Tests, alle BESTANDEN (\texttt{EinsteinRecoveryTests} als TEILWEISE
vermerkt).}

\chapter{Der Auditpfad der feindlichen Begutachtung}
\label{ch:audit-trail}

Dieses Kapitel ist die tabellarische Aufzeichnung der in
Kapitel~\ref{ch:review-history} erzählten Geschichte der feindlichen Begutachtung. Es
listet die ``FATAL''-Punkte aus Runde 2 und ihre endgültigen Klassifikationen auf.

Ein feindlicher Gutachter (Runde~2) warf zwölf FATAL-Punkte auf. Jeder wurde erneut gegen
die ausführbaren Tests bewertet. Die Klassifikationen sind:

\begin{center}
\begin{tabular}{p{6.5cm}l}
\toprule
FATAL-Problem aus Runde 2 & Klassifikation \\
\midrule
Primitive als mathematische Objekte undefiniert & TEILWEISE GELÖST \\
Schr\"odinger-Ableitung zirkulär/unvollständig & GELÖST \\
Eichableitungen = elementare Topologie & TEILWEISE GELÖST \\
Komplexitätsargument keine Ableitung & TEILWEISE GELÖST \\
Struktur/Inhalt-Trennung = Immunisierung & TEILWEISE GELÖST \\
Zusammengesetzte Objekte zugelassen & TEILWEISE GELÖST \\
Innere 3 ungelöst, während ``vollständig'' & GELÖST \\
T-09 bei 0,10 kann nicht schließen & GELÖST \\
``Strukturell vollständig''-Übertreibung & GELÖST \\
Gravitation$\to$Einstein ist ontologisch & GELÖST \\
$\mathrm{Aut}(S^1)=U(1)$ als EM-Ableitung & GELÖST \\
Kein kontrollierter Kontinuumsgrenzwert & TEILWEISE GELÖST \\
\bottomrule
\end{tabular}
\end{center}

\textbf{Bilanz: GELÖST = 6 · TEILWEISE GELÖST = 6 · OFFEN = 0.}

\section{Was sich änderte}

Die Rahmenübertreibungen (``geschlossen'' vs.\ ``disponiert'', ``No-Go'' vs.\
``bedingtes No-Go'', ``Ableitung'' vs.\ ``ontologische Reinterpretation'') wurden im
überarbeiteten Papier korrigiert. Der \emph{substantielle} Einwand — ``kein kontrollierter
Kontinuumsgrenzwert'' — wurde auf der Schr\"odinger-Seite beantwortet (nun getestet:
$L_Q\to-\nabla^2\to$ Schr\"odinger, sowie gekrümmtes $L_W$, Laplace--Beltrami, unitär)
und \emph{teilweise} auf der Einstein-Seite (die Standardkette funktioniert, aber die
Metrik und die BDG-Wirkung sind importiert).

\section{Der entscheidende Rest}

Der einzige wirklich offene Punkt ist die \textbf{native Metrik $\to$ Operator-Kopplung}
(die ``G4''-Lücke): AT importiert $g_{\mu\nu}$ (Malament), statt es aus Q-Ereignissen zu
erzeugen. Der Metrik\emph{ursprung} ist geschlossen
(Kapitel~\ref{ch:metric-emergence}); die Metrik\emph{dynamik} bleibt importiert.

\textbf{Urteil:} READY AS A FOUNDATION MONOGRAPH --- READY AS A RESEARCH-PROGRAM PAPER
--- NOT READY AS A DERIVATION PAPER (siehe Kapitel~\ref{ch:conclusion}).

%=============================================================================

\chapter{Was die feindliche Begutachtung überlebte}
\label{ch:survived}

Dieses Kapitel fasst die beiden Runden feindlicher Begutachtung in einer einzigen
Überlebenstabelle zusammen. Für jede Hauptbehauptung hält es den ursprünglichen Status,
die Kritik des Gutachters und den endgültigen Status nach der Neubewertung gegen die
abgeschlossenen Tests fest. Das Muster ist einheitlich: Nichts wurde zurückgezogen, aber
viele Übertreibungen wurden auf präzise, verteidigungsfähige Formen herabgestuft.

\begin{longtable}{p{3.2cm}p{2.4cm}p{3.0cm}p{5.8cm}}
\toprule
Behauptung & Ursprünglicher Status & Kritik der Begutachtung & Endgültiger Status \\
\midrule
\endhead
Primitive & ``vollständig formalisiert'' & als mathematische Objekte undefiniert & teilweise formalisiert (Maß/Wirkung als fehlend zugestanden) \\
Schr\"odinger-Gleichung & abgeleitet & zirkulär / unvollständig & abgeleitet und getestet ($L_Q\to-\nabla^2\to$ Schr\"odinger, unitär) \\
Eichgruppe & abgeleitet & elementare Topologie & $U(1)$ abgeleitet; $SU(2)$ emergent; $SU(3)$ kontingent (Struktur, nicht Dynamik) \\
Dimensionalität $3{+}1$ & Ableitung & keine Ableitung & Fensterschnitt (T-02 $=$ 0,85) \\
Struktur/Inhalt-Trennung & Organisationsprinzip & Immunisierung & beibehalten; Falsifizierbarkeit gezeigt (Neutrino-Koide) \\
``Strukturell vollständig'' & vollständig & Übertreibung & ``disponiert''; innere 3 als offen markiert \\
No-Go-Sätze & No-Go & zu stark & bedingtes No-Go (relativ zum getesteten Wegraum) \\
Gravitation $\to$ Einstein & Ableitung & ontologisch, keine Ableitung & ontologische Reinterpretation im schwachen Feld \\
$G=\ell^2c^3/\hbar$ & abgeleitet & Dimensionsanalyse & dimensionales / Einheitenkonsistenz-Ergebnis \\
Kontinuumsgrenzwert & kontrolliert & kein kontrollierter Grenzwert & Schr\"odinger-Seite kontrolliert und getestet; Einstein-Seite partiell \\
Metrikursprung & geschlossen & eine Lücke & Ursprung geschlossen (Malament-Import, bewiesen); Dynamik importiert (G4) \\
Eindeutige Vorhersage & testbar & keine diskriminierende Vorhersage & lebende Vorhersagen; noch keine eindeutig vs.\ SM$+\Lambda$CDM \\
Koide-Relation & reale Struktur & Numerologie & real, vorhersagekräftig, RG-stabil; Ursprung nicht ableitbar (T-08) \\
Eichanzahl $n=3$ & abgeleitet & nicht festgelegt & vorläufiges No-Go T-09 (0,10) \\
\bottomrule
\end{longtable}

\textbf{Lesart der Tabelle.} Jede ``abgeleitet''-Behauptung, die die feindliche
Begutachtung überlebte, wurde entweder ein getesteter Satz ($U(1)$, der
Schr\"odinger-Grenzwert) oder eine ausdrücklich umgrenzte Aussage (Fensterschnitt,
ontologische Reinterpretation, Einheitenkonsistenz). Jede von den Gutachtern behauptete
``Lücke'' löste sich in eine von zwei ehrlichen Kategorien auf: einen importierten
bewiesenen Satz (Malament) oder eine zugestandene Theorielimitierung (G4, die Eichanzahl,
die diskriminierende Vorhersage). Nichts in dieser Tabelle wurde zurückgezogen, und nichts
wurde still fallengelassen.

\chapterend
{Keine Behauptung wurde zurückgezogen; Übertreibungen wurden auf präzise Formen
herabgestuft; die einzige echte Theorielimitierung (G4) und die drei
Ableitungspapier-Blocker sind nun explizit.}
{Dieselben drei wissenschaftlichen Blocker (native Metrik-Kopplung, eindeutige
Vorhersage, native BDG-Wirkung), die nur den Ableitungspapier-Anspruch betreffen.}
{Aus den abgeschlossenen Tests und den beiden Runden feindlicher Begutachtung tabelliert;
auf die Tabelle selbst greift kein ausführbarer Test.}

\chapter{Schlussfolgerung}
\label{ch:conclusion}

THE Q-MODEL ist \emph{strukturell vollständig} in dem präzisen Sinn, dass jede im
Umfang liegende Frage \textbf{disponiert} ist — abgeleitet, nicht ableitbar,
unterbestimmt, kontingent oder als Berechnungsschicht akzeptiert.

Der einzige Rest echter wissenschaftlicher Spannung ist der \textbf{Knoten der inneren 3}
— warum die innere Multiplizität/Anzahl bei 3 sättigt —, disponiert als
ungelöst-kontingent mit einem vorläufigen Eichanzahl-No-Go (T-09, 0,10). Dies ist die
eine offene Tür der Theorie.

\textbf{Struktur ist ableitbar. Inhalt wird realisiert.}

\textbf{Publikationsstatus:} READY AS A FOUNDATION MONOGRAPH --- READY AS A
RESEARCH-PROGRAM PAPER --- NOT READY AS A DERIVATION PAPER. Der dynamische
Schr\"odinger-Kern ist ausführbar und getestet; die Einstein-Wiederherstellung bleibt
logisch, nicht mathematisch (importierte Metrik + importierte BDG-Wirkung, $G$
dimensional, keine eindeutige diskriminierende Vorhersage). Die
Grundlagenmonographie- und Forschungsprogramm-Ansprüche werden von den verbleibenden
Ableitungspapier-Blockern nicht berührt, weil eine Grundlagenmonographie legitimerweise
importierte bewiesene Sätze verwenden darf.

%=============================================================================
\appendix
%=============================================================================

\chapter{Vollständiges Testinventar}
\label{app:tests}

\begin{longtable}{lllc}
\toprule
Testdatei & Tests & Verifizierte Größe & Status \\
\midrule
\endhead
GraphLaplacianContinuumTests & 1 & $\lambda_k=(1/\Delta x^2)[2-2\cos(\pi k/(N{+}1))]\to(\pi k)^2$ & BESTANDEN \\
BDGOperatorContinuumTests & 1 & BDG$\to\square$, $O(h^2)$ & BESTANDEN \\
QuantumGravityBridgeTests & 3 & $L_Q$ PSD; BDG indefinit; inkompatibel & BESTANDEN \\
CurvedSpaceBridgeTests & 3 & kein $\Delta_g$; $\Delta_g\to\nabla^2$; fehlt & BESTANDEN \\
MetricOperatorTests & 4 & Kandidatenkonstruktionen & BESTANDEN \\
WeightedLaplacianTests & 4 & symmetrisch; nullsummig; PSD; ungewichteter Grenzfall & BESTANDEN \\
LaplaceBeltramiTests & 3 & flacher Grenzfall; Gewichte; S¹ ($2k{-}1$-Entartung) & BESTANDEN \\
CurvedSchrodingerTests & 3 & unitär; flacher Grenzfall; gekrümmter Operator & BESTANDEN \\
EinsteinRecoveryTests & 3 & TEILWEISE (kein vollständiges $G_{\mu\nu}$) & BESTANDEN \\
EinsteinTensorTests & 4 & flach/sphär. $\Gamma$, $K$, $R$, $G$ & BESTANDEN \\
EinsteinTensorIntegrationTests & 4 & Builder; 3-Sphäre $G=-g$ & BESTANDEN \\
MetricGenerationTests & 4 & Distanz VORHANDEN; Kandidat TEILWEISE & BESTANDEN \\
MetricEmergenceTests & 4 & Faktor $f=\rho^{2/d}$; $R=-0.6145$ & BESTANDEN \\
ConformalStructureTests & 3 & Axiome; Nullinvariante; Malament & BESTANDEN \\
MetricOriginTests & 3 & Klasse eindeutig; Faktor frei; importiert & BESTANDEN \\
\bottomrule
\end{longtable}

\chapter{Konfidenz- und Klassifikationstabellen}
\label{app:confidence}

Die Konfidenzzahlen sind unkalibrierte ordinale Ränge (Phase-149/156-Schätzungen), keine
A-posteriori-Wahrscheinlichkeiten. Sie ordnen, sie messen nicht.

\section{Konfidenztabelle}

\begin{center}
\begin{tabular}{llc}
\toprule
Ergebnis & Konfidenz & Anmerkung \\
\midrule
Programmkonsistenz & 0,95 & null Kategoriekonflikte \\
$U(1)$ & 0,95 & Satz \\
$N\ge3$ & 0,90 & CP-Untergrenze \\
Neutrino-Koide & 0,90 & Falsifikation (Berechnung) \\
räumliche 3 & 0,85 & Fensterschnitt \\
$SU(2)$ & 0,70 & emergent \\
Koide $45^\circ$ & 0,70 & strukturiert \\
$N\le3$ & 0,70 & No-Go \\
$SU(3)$-Struktur & 0,10 & vorläufiges No-Go \\
\bottomrule
\end{tabular}
\end{center}

\section{Endgültige Klassifikation}

\begin{center}
\begin{tabular}{ll}
\toprule
Frage & Klassifikation \\
\midrule
Ist der Metrikursprung geschlossen? & JA (importierter Malament, bewiesen) \\
Ist der Schr\"odinger-Kontinuumsgrenzwert kontrolliert? & JA (getestet) \\
Ist die Einstein-Wiederherstellung eine Ableitung? & NEIN (logisch, importiert) \\
Ist AT journalreif? & NEIN \\
Ist AT whitepaperreif? & JA \\
\bottomrule
\end{tabular}
\end{center}

%=============================================================================

\chapter{Ausführliche Ableitung: Die Einstein-Kette auf der 2-Sphäre}
\label{app:worked}

Dieser Anhang führt die übliche differentialgeometrische Kette explizit vor, damit die
Leserin oder der Leser jeden Schritt verfolgen kann, ohne den Repository-Code zu öffnen.

\section{Aufbau}

Die Einheits-2-Sphäre hat die Metrik (Koordinaten $x^1=\theta$, $x^2=\phi$)
\begin{equation}
g = \mathrm{diag}\!\left(g_{\theta\theta},\ g_{\phi\phi}\right)
  = \mathrm{diag}\!\left(1,\ \sin^2\theta\right),
\qquad g^{\theta\theta}=1,\quad g^{\phi\phi}=\frac{1}{\sin^2\theta}.
\end{equation}

\section{Christoffel-Symbole}

\begin{equation}
\Gamma^\lambda_{\mu\nu} = \tfrac12 g^{\lambda\sigma}\left(
\partial_\mu g_{\sigma\nu} + \partial_\nu g_{\sigma\mu} - \partial_\sigma g_{\mu\nu}\right).
\end{equation}

Die einzige nicht verschwindende partielle Ableitung ist
$\partial_\theta g_{\phi\phi} = 2\sin\theta\cos\theta$. Daher
\begin{align}
\Gamma^\theta_{\phi\phi} &= -\tfrac12\, g^{\theta\theta}\, \partial_\theta g_{\phi\phi}
= -\sin\theta\cos\theta, \\
\Gamma^\phi_{\theta\phi} = \Gamma^\phi_{\phi\theta}
&= \tfrac12\, g^{\phi\phi}\, \partial_\theta g_{\phi\phi}
= \frac{\cos\theta}{\sin\theta} = \cot\theta.
\end{align}
Bei $\theta=\pi/4$: $\Gamma^\theta_{\phi\phi}=-1/2$, $\Gamma^\phi_{\theta\phi}=1$.

\section{Riemann-Krümmung}

\begin{equation}
R^\rho_{\ \sigma\mu\nu} = \partial_\mu\Gamma^\rho_{\nu\sigma}
- \partial_\nu\Gamma^\rho_{\mu\sigma}
+ \Gamma^\rho_{\mu\lambda}\Gamma^\lambda_{\nu\sigma}
- \Gamma^\rho_{\nu\lambda}\Gamma^\lambda_{\mu\sigma}.
\end{equation}

Die einzige unabhängige Komponente ist
\begin{align}
R^\theta_{\ \phi\theta\phi}
&= \partial_\theta\Gamma^\theta_{\phi\phi} - \Gamma^\theta_{\phi\phi}\Gamma^\phi_{\theta\phi} \\
&= -\left(\cos^2\theta-\sin^2\theta\right) + \sin\theta\cos\theta\cdot\cot\theta \\
&= \sin^2\theta.
\end{align}
Bei $\theta=\pi/4$: $R^\theta_{\ \phi\theta\phi}=1/2$.

\section{Gauß-Krümmung, Ricci und Einstein}

Gauß-Krümmung:
\begin{equation}
K = \frac{R^\theta_{\ \phi\theta\phi}}{g_{\phi\phi}}
= \frac{\sin^2\theta}{\sin^2\theta} = 1.
\end{equation}

Ricci-Tensor (in 2D gilt $R_{\mu\nu}=K g_{\mu\nu}$):
\begin{equation}
R_{\theta\theta}=K g_{\theta\theta}=1,\qquad
R_{\phi\phi}=K g_{\phi\phi}=\sin^2\theta=\tfrac12 \text{ bei }\theta=\pi/4.
\end{equation}

Ricci-Skalar:
\begin{equation}
R = g^{\mu\nu}R_{\mu\nu} = 2K = 2.
\end{equation}

Einstein-Tensor (verschwindet identisch in 2D):
\begin{equation}
G_{\mu\nu} = R_{\mu\nu} - \tfrac12 R g_{\mu\nu} = 0.
\end{equation}

All dies wird von \texttt{EinsteinTensorBuilder} numerisch bis auf $10^{-3}$ reproduziert.
Ein nichttrivialer $G_{\mu\nu}$ erfordert Dimension $\ge3$; die Einheits-3-Sphäre liefert
$G_{\mu\nu}=-g_{\mu\nu}$ (Anhang~\ref{app:tests}).

%=============================================================================

\chapter{Forschungsprogramme}
\label{app:programs}

Das Ableitungsprogramm ist in Forschungsprogramme gegliedert, die jeweils durch
ausführbare Experimente abgeschlossen werden. Die vier Programme, die das technische
Gewicht dieser Monographie tragen --- das Eich-, Koide-, Kontinuums- und Metrikprogramm
---, werden in Kapitel~\ref{ch:journey} erzählt; dieser Anhang hält das
Programmregister fest.

\begin{center}
\begin{tabular}{llc}
\toprule
Programm & Umfang & Experimente \\
\midrule
DATA & Kosmologie \& RAR & 10 (DATA-001$\to$010) \\
QM & Quantengrundlagen & 5 (QM-001$\to$005) \\
QG & Quantengravitation & 31 (QG-001$\to$031) \\
XC & Kontinuumsgrenzwert-Brücke & (Audits + Tests) \\
\bottomrule
\end{tabular}
\end{center}

Kernergebnisse des gesamten Programms:
\begin{itemize}
\item $G$ ist nicht fundamental: $G=\ell^2c^3/\hbar$ (Gravitation entsteht aus dem
Tripel).
\item $(\ell,\tau,\hbar)$ ist das irreduzible Tripel — ein Prozess, drei Aspekte.
\item Oszillation ist eine logische Unvermeidlichkeit (nicht entfernbar, ohne $Q$ zu
zerstören).
\item Gravitation ist ein Phasengradienten-Phänomen; Anziehung dominiert, Abstoßung ist
dunkle Energie.
\item Gravitationsmanipulation ist nicht möglich (QG-023$\to$031).
\item RAR-Ableitung $g_\dagger=cH_0/(2\pi)$ mit null freien Parametern.
\item Parameterkompression SM+GR ($\sim$26) $\to$ AT ($2{+}3$), $\sim$5--8$\times$
Reduktion.
\item Neue Vorhersagen: $w(z)=-1+0.015(1+z)^{3/2}$; evolvierendes $g_\dagger(z)$;
Euclid-Empfindlichkeit.
\end{itemize}

\chapter{Entwicklungsregister (Auswahl)}
\label{app:ledger}

Dieser Anhang hält in kompakter Form eine Auswahl der Experimente und Audits fest, die
die Erzählung von Kapitel~\ref{ch:timeline} untermauern. Er ist nicht erschöpfend: Das
vollständige Laborbuch liegt in der Projektdokumentation. Jeder Eintrag nennt den
Bezeichner, die Frage oder Behauptung und das festgehaltene Ergebnis. Die Einträge sind
nach Ära gruppiert.

\begin{longtable}{p{4.2cm}p{5.2cm}p{5.6cm}}
\toprule
Bezeichner & Frage / Behauptung & Ergebnis / Disposition \\
\midrule
\endhead
\multicolumn{3}{l}{\emph{Frühe Oszillatorexperimente}} \\
\midrule
AT-001 & Entsteht Synchronisation? & Ja: Ordnungsparameter $R\to1$. \\
AT-002 & Ergeben Zufallsmatrizen Emergenz? & Nein: Wigner-Spektrum, keine strukturierte Emergenz. \\
AT-003 & Ist statische Topologie ausreichend? & Nein: Dynamik überwiegt statische Struktur. \\
AT-004 & Synchronisiert ein Feld zwei asymmetrische Oszillatoren? & Nein: Feld wirkt als Resonanzmedium. \\
AT-005 & Bilden sich Resonanzcluster? & Ja: zwei Cluster aus drei Oszillatoren, $\sim5\%$ Energiebeteiligung. \\
AT-006 & Gibt es eine kritische Resonanzdichte? & Ja: $\rho_c\approx0.09$, perkolationsartiger Übergang. \\
AT-007 & Wie viele Resonanzfamilien existieren? & Fünf gemeldet. \\
AT-008 & Welche Familien sind reproduzierbar? & Nur F4 (74\%, Score 0,721); der Rest sind Keim-Artefakte. \\
\midrule
\multicolumn{3}{l}{\emph{Alternative Grundlagen (ResearchX)}} \\
\midrule
X001 & Audit verborgener Annahmen & 11 gefunden; der statische Graph ist Nr.~1. \\
X002--X004 & Innovieren dynamische Graphen? & Nein: Knotenbewegung, Mobilität, Wachstum fügen nichts Offenes hinzu. \\
X005--X006 & Erhält Nichtlinearität die Quanten-Lesart? & Nein: Eigenmoden zerbrechen in Solitonen; Artefakte der linearen Algebra. \\
X007--X008 & Universelle Artentheorie & 16 Trägerklassen über fünf Regime. \\
X009--X015 & Tiefste Grundlagen & Reversibilität $+$ Selbstkonsistenz ist minimal hinreichend; Rev$\cap$SC $=$ QM. \\
X016--X017 & Realitätskarte & $(R,S)$ ist eine Karte, kein dynamischer Fluss. \\
X018--X022 & Komplexitätstreppe & Stufe 6 nicht beobachtet; Operatorevolution nötig. \\
X023--X027 & Suche nach offener Evolution & Verzögerte Sättigung; Schubfachbeweis für endliche Systeme. \\
X028--X029 & Konsequenzen endlicher Universen & Sättigung physikalisch irrelevant; hybride Architekturen optimal. \\
X030--X031 & Quantenoptimalität & QM ist das eindeutige globale Maximum endlicher Komplexität. \\
X032--X034 & Vereinheitlichung & $L_Q$ ist nicht fundamental; Emergenzlücke geschlossen. \\
X035--X036 & Ursprung von $Q$; Komplexitätssatz & $Q$ irreduzibel; Schr\"odinger abgeleitet (Klassifikation C). \\
X037--X039 & Born, Messung, Zufall & Born abgeleitet ($\alpha=2$); Viele-Welten inkompatibel; Zufall irreduzibel. \\
X040--X042 & Zeit, Gravitation, Dimensionalität & Zeit $=$ Halbordnung; Gravitation $=$ Kausalmenge; $3{+}1$ abgeleitet. \\
X043--X046 & Fundamentale Skalen & $G=\ell^2c^3/\hbar$; $\hbar=1$; $c,G,\hbar$ vereinheitlicht; $\Lambda\sim H^2$. \\
X047--X049 & Materie und Eichstruktur & Teilchen $=$ Defekte; Eichstruktur entsteht; Selektionslücke. \\
X051--X056, X059--X060 & Teilchenphysik & Generationen, Massen, Mischung, Neutrinos; SM-Gruppe bevorzugt, nicht abgeleitet. \\
X057--X058, X060b--X060f & Parameterkompression & Zwei Primitive $+$ $M^2$; $M^2$ irreduzibel. \\
\midrule
\multicolumn{3}{l}{\emph{Späte Phasen (140--159)}} \\
\midrule
Phase 140 & Kausalmengen-$\Lambda$ & $\Lambda\sim1/\sqrt N$ (Postdiktion). \\
Phase 144 & Radialbeschleunigungsrelation & $g_\dagger=cH_0/2\pi$; das $2\pi$ ist ein numerischer Zufall. \\
Phase 148 & Flavour-/Yukawa-Reduzierbarkeit & Koide $Q=2/3$ kontingent, nicht abgeleitet. \\
Phase 149 & Eichursprung & $U(1)$ abgeleitet; $SU(2)$ emergent; $SU(3)$ kontingent. \\
Phase 150 & Multiplizität ``3'' & Raumzeit-3 abgeleitet; innere 3 selektiert. \\
Phase 151 & Obergrenze $N\le3$ & Empirisch/kontingent; keine Stabilitäts-/Anomalie-/Darstellungsschranke. \\
Phase 152 & Ensembles der Zufälligen Aktualisierung & Vier unabhängige Ensembles unter einem Lognormalgesetz. \\
Phase 153 & Kaskadenvereinheitlichung & Drei Universalitätsklassen sind unabhängige Kaskaden. \\
Phase 155 & Neutrino-Koide & Falsifiziert ($Q_{\max}=0.585/0.500<2/3$). \\
Phase 159 & Koide-Abschluss & Sieben Wege, sechs No-Gos, einer blockiert; geschlossene Frage. \\
QG-080--100 & Kausalmengen-Grundlagenprogramm & $\Lambda\sim H^2/c^2$; $a_0\sim cH$; Endantwort erreicht. \\
Feindliche Begutachtung (Runde~1) & v1.0-Papier-Audit & 1 gültig, 11 teilweise gültig, 0 ungültig; Urteil NOT READY. \\
Feindliche Begutachtung (Runde~2) & FATAL-Punkte-Audit & 6 gelöst, 6 teilweise gelöst, 0 offen. \\
v1.0-Veröffentlichung & Versionsreife & Vier Dispositionen; v1.0 veröffentlicht. \\
\bottomrule
\end{longtable}

\chapter{Glossar}
\label{app:glossary}

\begin{center}
\begin{tabular}{lp{10cm}}
\toprule
Begriff & Bedeutung \\
\midrule
$Q$ & Individuationsprinzip (Ontologie) \\
Zufällige Aktualisierung & ontologischer Zufall (Annahme A-03) \\
$(\ell,\tau,\hbar)$ & Raumzeit-Skala, Uhr, Wirkung (Einheiten) \\
$M^2$ & Nichtlinearitätsparameter (einziger kontinuierlicher) \\
$L_Q$ & Graphen-Laplace-Operator $D-A$ \\
$L_W$ & gewichteter Laplace-Operator $D_K-K$ \\
BDG & Benincasa--Dowker--Glaser-d'Alembert-Operator \\
$\Delta_g$ & Laplace--Beltrami-Operator \\
$G_{\mu\nu}$ & Einstein-Tensor $R_{\mu\nu}-\tfrac12 R g_{\mu\nu}$ \\
\derived & aus Primitiven per Satz berechenbar \\
\realunder & reale Struktur, nicht ableitbarer Ursprung \\
\drawn & zufällige Ziehung des Häufigkeitsgesetzes \\
T-08\dots T-12 & bedingte No-Go-Sätze \\
A-03, A-06, A-07, A-10 & Annahmen (zufällige Aktualisierung; $Z_2$-Anhebung; Eichdynamik; $N$/$n$-Trennung) \\
\bottomrule
\end{tabular}
\end{center}

%=============================================================================

\chapter{Literatur}
\label{app:references}

\begin{thebibliography}{99}

\bibitem{malament1977}
D.~B.~Malament, ``The class of continuous timelike curves determines the topology of
spacetime,'' \emph{J. Math. Phys.} \textbf{18}, 1399 (1977).

\bibitem{hkm1976}
S.~W.~Hawking, A.~R.~King, P.~J.~McCarthy, ``A new topology for curved space--time which
incorporates the causal, differential, and conformal structures,'' \emph{J. Math. Phys.}
\textbf{17}, 174 (1976).

\bibitem{myrheim1978}
J.~Myrheim, ``Statistical geometry,'' CERN-TH-2538 (1978).

\bibitem{bdg2010}
D.~M.~T.~Benincasa, F.~Dowker, ``Scalar curvature of a causal set,'' \emph{Phys. Rev. Lett.}
\textbf{104}, 181301 (2010); L.~Glaser, \emph{Class. Quant. Grav.} \textbf{31}, 095007
(2014).

\bibitem{gleason1957}
A.~M.~Gleason, ``Measures on the closed subspaces of a Hilbert space,''
\emph{J. Math. Mech.} \textbf{6}, 885 (1957).

\bibitem{sorkin}
R.~D.~Sorkin, ``Causal sets: discrete gravity,'' \emph{Proceedings of the Valdivia Summer
School} (2003).

\bibitem{belkin2003}
M.~Belkin, P.~Niyogi, ``Laplacian eigenmaps for dimensionality reduction and data
representation,'' \emph{Neural Computation} \textbf{15}, 1373 (2003).

\bibitem{koide1981}
Y.~Koide, ``Fermion mass relation and the generation structure,'' \emph{Phys. Rev. D}
\textbf{28}, 252 (1983) [letter 1981].

\bibitem{bertrand}
J.~Bertrand, ``Th\'eor\`eme relatif au mouvement d'un point attir\'e vers un centre fixe,''
\emph{C. R. Acad. Sci.} \textbf{77}, 849 (1873).

\end{thebibliography}

\end{document}
